%%%%%%%%%%%%%%%%%%%%%%%%%%%%%%%%%%%%%%%%%%%%%%%%%%%%%%%%%%%%%%%%%%%%%%%%%%%%%%%%
%% Autonomous Navigation Through Categorical State Counting
%% in Coupled Oscillator Networks
%%
%% Paper II of III: Autonomous Driving via Trajectory Completion Computing
%%
%% Author: Kundai Farai Sachikonye
%% Technical University of Munich
%%%%%%%%%%%%%%%%%%%%%%%%%%%%%%%%%%%%%%%%%%%%%%%%%%%%%%%%%%%%%%%%%%%%%%%%%%%%%%%%

\documentclass[twocolumn,10pt]{article}

\usepackage{amsmath}
\usepackage{amssymb}
\usepackage{amsthm}
\usepackage{graphicx}
\usepackage{hyperref}
\usepackage{booktabs}
\usepackage{siunitx}
\usepackage{geometry}
\usepackage{algorithm}
\usepackage{algorithmic}
\usepackage{mathtools}
\usepackage{bm}
\usepackage{enumitem}

\geometry{margin=1in}

\hypersetup{
  colorlinks=true,
  linkcolor=blue,
  citecolor=blue,
  urlcolor=blue
}

% Theorem environments
\newtheorem{theorem}{Theorem}[section]
\newtheorem{lemma}[theorem]{Lemma}
\newtheorem{proposition}[theorem]{Proposition}
\newtheorem{corollary}[theorem]{Corollary}
\newtheorem{axiom}[theorem]{Axiom}
\theoremstyle{definition}
\newtheorem{definition}[theorem]{Definition}
\theoremstyle{remark}
\newtheorem{remark}[theorem]{Remark}

% Custom commands
\newcommand{\Sosc}{S_{\mathrm{osc}}}
\newcommand{\Scat}{S_{\mathrm{cat}}}
\newcommand{\Spart}{S_{\mathrm{part}}}
\newcommand{\Sent}{\mathbf{S}}
\newcommand{\kB}{k_{\mathrm{B}}}
\newcommand{\Ocat}{\hat{O}_{\mathrm{cat}}}
\newcommand{\Ophys}{\hat{O}_{\mathrm{phys}}}
\newcommand{\R}{\mathbb{R}}
\newcommand{\N}{\mathbb{N}}
\newcommand{\Z}{\mathbb{Z}}
\newcommand{\Ocal}{\mathcal{O}}
\newcommand{\Tcal}{\mathcal{T}}
\newcommand{\Gcal}{\mathcal{G}}
\newcommand{\Hcal}{\mathcal{H}}
\newcommand{\Dcal}{\mathcal{D}}
\newcommand{\Fcal}{\mathcal{F}}
\newcommand{\Ccal}{\mathcal{C}}
\newcommand{\Mcal}{\mathcal{M}}
\newcommand{\tauP}{\tau_{\mathrm{p}}}

\title{Autonomous Navigation Through Categorical State Counting\\
in Coupled Oscillator Networks}

\author{Kundai Farai Sachikonye\\
Technical University of Munich\\
\texttt{kundai.sachikonye@wzw.tum.de}}

\date{March 2026}

\begin{document}

\maketitle

\begin{abstract}
We present a rigorous mathematical framework for autonomous vehicle navigation
that replaces forward simulation and trajectory prediction with categorical
state counting in coupled oscillator networks. The central result is the
\emph{Triple Process Identity}: observation $O(x)$, computing $C(x)$, and
processing $P(x)$ are categorically equivalent when realized on physical
oscillators occupying bounded phase space. A vehicle's constituent oscillators
--- CPU clock ($\sim$\SI{}{GHz}), engine crankshaft ($\sim$\SI{50}{Hz}),
wheel rotation ($\sim$\SI{10}{Hz}), and atmospheric molecular vibrations
($\sim$\SI{e13}{Hz}) --- form a harmonic coincidence network whose
phase-locked rational frequency ratios define a counting loop. We prove that
this counting loop computes navigation functions with complexity
$\Ocal(\log_3 M)$ for state finding, $\Ocal(n^k)$ for sufficiency checking,
and $\Ocal(1)$ for recognition, where $M$ is the partition count. Through the
Position-Partition Bijection $\Pi: \R^3 \to [0,1]^3$, we show that
atmospheric S-entropy fields encode absolute position without GPS. The
oscillator--processor duality $\omega \equiv R_{\mathrm{compute}}$ and the
commutation relation $[\Ocat, \Ophys] = 0$ ensure that categorical
observation and physical measurement yield identical navigation states. We
derive trans-Planckian timing enhancement of $10^{120.95}$ effective
resolution through five vehicle-specific mechanisms and prove that other
vehicles are detected as perturbations of the local S-entropy field,
requiring no predictive model. The framework is grounded in Poincar\'e
recurrence, Boltzmann statistical mechanics, and Kuramoto synchronization
theory. We propose an experimental protocol using instrumented vehicles to
validate the theory against conventional GPS-based navigation.

\medskip
\noindent\textbf{Keywords:} autonomous navigation, oscillator networks, categorical state counting,
Poincar\'e recurrence, S-entropy, harmonic coincidence, trajectory completion,
coupled oscillators, trans-Planckian timing, sufficiency recognition
\end{abstract}

%%%%%%%%%%%%%%%%%%%%%%%%%%%%%%%%%%%%%%%%%%%%%%%%%%%%%%%%%%%%%%%%%%%%%%%%%%%%%%%%
\section{Introduction}\label{sec:intro}
%%%%%%%%%%%%%%%%%%%%%%%%%%%%%%%%%%%%%%%%%%%%%%%%%%%%%%%%%%%%%%%%%%%%%%%%%%%%%%%%

The prevailing paradigm in autonomous driving research treats navigation as a
prediction problem: given the current state of the world, forecast future
trajectories and select the optimal path~\cite{thrun2006,bojarski2016}. This
approach inherits the computational burden of forward simulation---solving
differential equations, propagating uncertainty, and evaluating cost
functions over exponentially branching futures. We argue that this paradigm
rests on a fundamental conceptual error that we term the \emph{prediction fallacy}.

\begin{definition}[Prediction Fallacy]\label{def:prediction-fallacy}
The prediction fallacy is the assumption that navigation requires a forward
model $f: \mathcal{S}_t \to \mathcal{S}_{t+\Delta t}$ mapping current states
to future states. Formally, it is the claim that safe navigation necessitates
computing
\begin{equation}\label{eq:prediction}
  s(t + \Delta t) = f(s(t), u(t), w(t))
\end{equation}
where $s(t) \in \mathcal{S}$ is the state, $u(t)$ the control input, and
$w(t)$ the environment process, for some horizon $\Delta t > 0$.
\end{definition}

Human drivers refute this assumption at every moment. An experienced driver
does not solve differential equations or propagate probability distributions
over possible futures. Instead, the driver \emph{recognizes sufficiency}:
the current configuration of sensory inputs either suffices for continued
safe operation or it does not. This recognition is instantaneous, requiring
no simulation.

The present paper---the second in a three-paper series on autonomous driving
via trajectory completion computing~\cite{sachikonye2026trajectory}---formalizes
this observation through the theory of categorical state counting in coupled
oscillator networks. The key mathematical structure is:

\begin{axiom}[Triple Process Identity]\label{ax:triple-process}
For any physical system $\mathcal{P}$ occupying bounded phase space with
oscillatory dynamics at frequency $\omega$, the observation operator
$O(x)$, the computation operator $C(x)$, and the processing operator
$P(x)$ satisfy
\begin{equation}\label{eq:triple-identity}
  O(x) \equiv C(x) \equiv P(x).
\end{equation}
That is, observation IS computing IS processing. These are not merely
correlated or approximately equal; they are categorically identical morphisms
in the category of physical processes.
\end{axiom}

The consequences for autonomous navigation are immediate. If observation is
computing, then a vehicle's physical oscillators---the CPU clock, the engine
crankshaft, the wheel rotations, the atmospheric molecular vibrations---are
not merely producing data for a separate computational system. They \emph{are}
the computational system. Navigation reduces to counting the categorical
states that these coupled oscillators traverse.

\begin{definition}[Oscillator--Processor Duality]\label{def:osc-proc-duality}
For any oscillator with angular frequency $\omega$ and any computational
process with rate $R_{\mathrm{compute}}$ (operations per second), the
oscillator--processor duality is the identity
\begin{equation}\label{eq:osc-proc}
  \omega \equiv R_{\mathrm{compute}}.
\end{equation}
Each oscillation cycle constitutes exactly one computational operation: the
transition from one categorical state to the next in the bounded phase space.
\end{definition}

This duality was established in~\cite{sachikonye2026poincare} through the
chain: bounded phase space $\Rightarrow$ Poincar\'e
recurrence~\cite{poincare1890} $\Rightarrow$ oscillation $\Rightarrow$
categorical states $\Rightarrow$ counting $\Rightarrow$ computation. The
present paper applies this chain specifically to the vehicle navigation
problem.

The paper is organized as follows. Section~\ref{sec:vehicle-oscillator}
establishes the vehicle as a multi-scale oscillator network and proves the
fundamental counting identity. Section~\ref{sec:harmonic-coincidence}
develops the harmonic coincidence network through which rational frequency
ratios enable phase-locked navigation. Section~\ref{sec:spectrometer} shows
how the vehicle's computer functions as a spectrometer via
precision-by-difference. Section~\ref{sec:multimodal} constructs the
multi-modal observer architecture. Section~\ref{sec:sufficiency} formalizes
sufficiency recognition as the alternative to prediction.
Section~\ref{sec:transplanckian} derives the trans-Planckian timing
enhancement. Section~\ref{sec:position} proves position determination
without GPS via the Position-Partition Bijection.
Section~\ref{sec:other-vehicles} treats other vehicles as categorical state
perturbations. Section~\ref{sec:experiment} proposes experimental validation.
Section~\ref{sec:conclusion} concludes.


%%%%%%%%%%%%%%%%%%%%%%%%%%%%%%%%%%%%%%%%%%%%%%%%%%%%%%%%%%%%%%%%%%%%%%%%%%%%%%%%
\section{Vehicle as Oscillator Network}\label{sec:vehicle-oscillator}
%%%%%%%%%%%%%%%%%%%%%%%%%%%%%%%%%%%%%%%%%%%%%%%%%%%%%%%%%%%%%%%%%%%%%%%%%%%%%%%%

\subsection{Oscillator Inventory}

A modern vehicle contains oscillators spanning over thirteen orders of
magnitude in frequency. We catalog the principal oscillators relevant to
navigation.

\begin{definition}[Vehicle Oscillator Inventory]\label{def:osc-inventory}
A vehicle oscillator inventory is a set
$\mathcal{V} = \{(\omega_i, \Phi_i, \mathcal{H}_i)\}_{i=1}^{N}$
where $\omega_i$ is the angular frequency, $\Phi_i$ the bounded phase
space, and $\mathcal{H}_i$ the Hilbert space of categorical states for the
$i$-th oscillator. The principal oscillators are:
\begin{enumerate}[label=(\roman*)]
  \item \textbf{CPU clock}: $\omega_{\mathrm{CPU}} \sim 2\pi \times
    \SI{3e9}{Hz}$, $\Phi_{\mathrm{CPU}} = S^1$ (ring oscillator phase),
    $\dim \mathcal{H}_{\mathrm{CPU}} = 2^{64}$.
  \item \textbf{Sensor clocks}: $\omega_{\mathrm{sensor}} \sim 2\pi \times
    \SI{e6}{Hz}$ to $2\pi \times \SI{e9}{Hz}$ (LiDAR, radar, camera pixel
    clocks).
  \item \textbf{Engine crankshaft}: $\omega_{\mathrm{engine}} \sim 2\pi
    \times \SI{50}{Hz}$ at \SI{3000}{rpm}.
  \item \textbf{Wheel rotation}: $\omega_{\mathrm{wheel}} \sim 2\pi \times
    \SI{10}{Hz}$ at \SI{100}{km/h} with standard tires.
  \item \textbf{Suspension oscillation}: $\omega_{\mathrm{susp}} \sim 2\pi
    \times \SI{1.5}{Hz}$ (typical body resonance).
  \item \textbf{Atmospheric molecular vibration}: $\omega_{\mathrm{atm}}
    \sim 2\pi \times \SI{e13}{Hz}$ (rotational--vibrational modes of
    N$_2$, O$_2$, CO$_2$)~\cite{herzberg1945}.
\end{enumerate}
\end{definition}

The frequency span from suspension (\SI{1.5}{Hz}) to atmospheric vibration
(\SI{e13}{Hz}) covers approximately 13 decades. This span is not incidental;
it is the basis for the trans-Planckian timing enhancement derived in
Section~\ref{sec:transplanckian}.

\begin{figure}[t]
  \centering
  % \includegraphics[width=\columnwidth]{figures/oscillator_inventory.pdf}
  \rule{\columnwidth}{4cm}
  \caption{Vehicle oscillator inventory spanning 13 orders of magnitude.
  Horizontal axis: frequency (Hz) on logarithmic scale from $10^0$ to
  $10^{13}$. Vertical bars indicate the principal oscillators: suspension
  (\SI{1.5}{Hz}), wheel (\SI{10}{Hz}), engine (\SI{50}{Hz}), sensor clocks
  (\SI{e6}{Hz}--\SI{e9}{Hz}), CPU (\SI{3e9}{Hz}), and atmospheric
  molecular modes (\SI{e13}{Hz}). The 13-decade span enables trans-Planckian
  enhancement through harmonic coincidence.}
  \label{fig:oscillator-inventory}
\end{figure}

\subsection{Counting Loop Definition}

\begin{definition}[Counting Loop]\label{def:counting-loop}
Let $\mathcal{V} = \{(\omega_i, \Phi_i, \mathcal{H}_i)\}_{i=1}^{N}$ be a
vehicle oscillator inventory. A \emph{counting loop} is a closed cycle
$\gamma$ in the product phase space
$\Phi = \prod_{i=1}^N \Phi_i$
such that each oscillator $i$ completes an integer number $n_i \in \N$ of
full cycles:
\begin{equation}\label{eq:counting-loop}
  \gamma: [0, T_\gamma] \to \Phi, \qquad
  \omega_i \, T_\gamma = 2\pi n_i \quad \forall\, i \in \{1, \ldots, N\}.
\end{equation}
The counting loop period is the least common period:
\begin{equation}\label{eq:loop-period}
  T_\gamma = \mathrm{lcm}\!\left(\frac{2\pi}{\omega_1}, \ldots,
  \frac{2\pi}{\omega_N}\right).
\end{equation}
The total count accumulated in one loop is
\begin{equation}\label{eq:total-count}
  \mathcal{N}_\gamma = \sum_{i=1}^{N} n_i = \sum_{i=1}^{N}
  \frac{\omega_i \, T_\gamma}{2\pi}.
\end{equation}
\end{definition}

\begin{remark}\label{rem:rational-approx}
In practice, the frequency ratios $\omega_i / \omega_j$ are irrational, so
exact counting loops do not close. However, for any $\epsilon > 0$, there
exist $\epsilon$-approximate counting loops with $|\omega_i T_\gamma -
2\pi n_i| < \epsilon$ for all $i$, by the multidimensional equidistribution
theorem~\cite{strogatz2000}. The deviation from exact closure is precisely
the information that encodes position, as shown in
Section~\ref{sec:position}.
\end{remark}

\subsection{The Fundamental Counting Identity}

We now prove the central identity connecting oscillator entropy, categorical
entropy, and partition entropy.

\begin{theorem}[Triple Entropy Equivalence]\label{thm:triple-entropy}
Let $\mathcal{V}$ be a vehicle oscillator inventory with $M$ distinguishable
states across all oscillators and $n$ partition levels. Then the oscillator
entropy $\Sosc$, the categorical entropy $\Scat$, and the partition entropy
$\Spart$ satisfy
\begin{equation}\label{eq:triple-entropy}
  \Sosc = \Scat = \Spart = \kB M \ln n.
\end{equation}
\end{theorem}

\begin{proof}
We establish this through two equalities.

\emph{Step 1: $\Sosc = \Scat$.} Each oscillator $(\omega_i, \Phi_i,
\mathcal{H}_i)$ occupies a bounded phase space $\Phi_i$. By the Poincar\'e
recurrence theorem~\cite{poincare1890}, every trajectory in $\Phi_i$
returns arbitrarily close to its initial point. This recurrence generates a
discrete sequence of return times $\{t_k^{(i)}\}_{k=1}^{\infty}$. Each
return constitutes a transition between categorical states in $\mathcal{H}_i$.
The number of distinguishable categorical states accessible to oscillator
$i$ is $|\mathcal{H}_i|$, so its categorical entropy
is $\kB \ln |\mathcal{H}_i|$~\cite{boltzmann1896}.

The total oscillator entropy is
\begin{equation}
  \Sosc = \kB \sum_{i=1}^{N} \ln |\mathcal{H}_i|
  = \kB \ln \prod_{i=1}^{N} |\mathcal{H}_i|.
\end{equation}
The total categorical state space is $\mathcal{H} = \bigotimes_{i=1}^N
\mathcal{H}_i$ with $|\mathcal{H}| = \prod_i |\mathcal{H}_i|$, so
$\Scat = \kB \ln |\mathcal{H}| = \Sosc$.

\emph{Step 2: $\Scat = \Spart$.} The partition coordinates
$(n, \ell, m, s)$~\cite{sachikonye2026partition} assign to each categorical
state a unique address in a hierarchical partition with degeneracy
$C(n) = 2n^2$~\cite{sachikonye2026poincare}. For $M$ states distributed
across $n$ levels, the partition entropy is
\begin{equation}
  \Spart = \kB \sum_{j=1}^{n} C(j) \ln C(j)
\end{equation}
which, by the partition--category
bijection~\cite{sachikonye2026partition}, equals $\kB M \ln n$ when the
states are maximally distributed. By the maximum entropy
principle~\cite{jaynes1957}, the physical distribution is precisely this
maximum, establishing $\Scat = \Spart = \kB M \ln n$.
\end{proof}

\begin{corollary}[Counting Rate Identity]\label{cor:counting-rate}
The rate of state counting equals the mean reciprocal Poincar\'e return time:
\begin{equation}\label{eq:counting-rate}
  \frac{dM}{dt} = \frac{\omega}{2\pi / M} = \frac{1}{\langle \tauP \rangle}
\end{equation}
where $\langle \tauP \rangle$ is the mean Poincar\'e return time averaged
over all oscillators in $\mathcal{V}$.
\end{corollary}

\begin{proof}
Each oscillator $i$ with frequency $\omega_i$ completes $\omega_i / (2\pi)$
cycles per unit time. Each cycle increments the state count by one. The
total rate across the network is
\begin{equation}
  \frac{dM}{dt} = \sum_{i=1}^N \frac{\omega_i}{2\pi}.
\end{equation}
The mean Poincar\'e return time is $\langle \tauP \rangle = N / \sum_i
(\omega_i / 2\pi)$, from which the identity follows for the network-averaged
frequency $\omega = (1/N)\sum_i \omega_i$ with $M = N$.
\end{proof}

\begin{definition}[S-Entropy Vector]\label{def:s-entropy}
The \emph{S-entropy} of a vehicle oscillator network is the vector
\begin{equation}\label{eq:s-entropy}
  \Sent = (S_k, S_t, S_e) \in [0,1]^3
\end{equation}
where:
\begin{itemize}
  \item $S_k \in [0,1]$ is the \emph{kinematic S-entropy}: the normalized
    entropy of the wheel and engine oscillator subsystem, encoding velocity
    and acceleration.
  \item $S_t \in [0,1]$ is the \emph{temporal S-entropy}: the normalized
    entropy of the clock oscillator subsystem, encoding timing precision.
  \item $S_e \in [0,1]$ is the \emph{environmental S-entropy}: the normalized
    entropy of the atmospheric and sensor oscillator subsystem, encoding
    environmental state.
\end{itemize}
Each component is normalized by its respective maximum:
$S_\alpha = \Sosc^{(\alpha)} / \Sosc^{(\alpha),\max}$ for
$\alpha \in \{k, t, e\}$.
\end{definition}

\begin{theorem}[S-Entropy Completeness]\label{thm:s-completeness}
The S-entropy vector $\Sent = (S_k, S_t, S_e)$ provides a complete
description of the vehicle's navigational state. That is, two
navigational configurations are identical if and only if they have the same
S-entropy vector.
\end{theorem}

\begin{proof}
Let $\mathcal{S}$ denote the space of navigational states and let
$\sigma: \mathcal{S} \to [0,1]^3$ be the S-entropy map. We must show
$\sigma$ is injective.

Suppose $\sigma(s_1) = \sigma(s_2)$ for $s_1, s_2 \in \mathcal{S}$.
Then $S_k(s_1) = S_k(s_2)$, which by the partition--category bijection
for the kinematic subsystem implies identical velocity--acceleration states.
Similarly, $S_t(s_1) = S_t(s_2)$ implies identical timing states, and
$S_e(s_1) = S_e(s_2)$ implies identical environmental states. Since the
navigational state is fully determined by kinematics, timing, and
environment (Axiom~\ref{ax:triple-process}), we have $s_1 = s_2$.

Surjectivity onto the image $\sigma(\mathcal{S}) \subseteq [0,1]^3$
follows from the fact that each S-entropy component ranges continuously
over $[0,1]$ as the physical parameters vary over their bounded ranges.
\end{proof}


%%%%%%%%%%%%%%%%%%%%%%%%%%%%%%%%%%%%%%%%%%%%%%%%%%%%%%%%%%%%%%%%%%%%%%%%%%%%%%%%
\section{Harmonic Coincidence Network}\label{sec:harmonic-coincidence}
%%%%%%%%%%%%%%%%%%%%%%%%%%%%%%%%%%%%%%%%%%%%%%%%%%%%%%%%%%%%%%%%%%%%%%%%%%%%%%%%

\subsection{Rational Frequency Ratios and Phase Locking}

The oscillators in $\mathcal{V}$ do not operate independently. Physical
coupling (mechanical, electromagnetic, acoustic) creates frequency
relationships that the counting loop exploits.

\begin{definition}[Harmonic Coincidence Graph]\label{def:harmonic-graph}
The \emph{harmonic coincidence graph} is the weighted graph
$\Gcal = (V, E, w)$ where:
\begin{itemize}
  \item $V = \{v_i\}_{i=1}^N$ corresponds to the oscillators in
    $\mathcal{V}$.
  \item $(v_i, v_j) \in E$ if and only if $\omega_i / \omega_j \approx
    p/q$ for some $p, q \in \N$ with $p + q \leq K$ (the
    \emph{coincidence threshold}).
  \item The weight $w(v_i, v_j) = 1/(p + q)$ measures the strength of the
    rational approximation.
\end{itemize}
\end{definition}

\begin{figure}[t]
  \centering
  % \includegraphics[width=\columnwidth]{figures/harmonic_graph.pdf}
  \rule{\columnwidth}{5cm}
  \caption{Harmonic coincidence graph for a vehicle oscillator network.
  Nodes represent principal oscillators (labels indicate frequency). Edges
  connect oscillators with rational frequency ratios $p/q$ with $p+q \leq 20$.
  Edge thickness is proportional to the weight $w = 1/(p+q)$. The
  atmospheric node (\SI{e13}{Hz}) connects to the CPU node (\SI{3e9}{Hz})
  through the approximate ratio $3333:1$. The engine (\SI{50}{Hz}) connects
  to the wheel (\SI{10}{Hz}) through the exact ratio $5:1$.}
  \label{fig:harmonic-graph}
\end{figure}

\begin{theorem}[Phase-Lock Existence]\label{thm:phase-lock}
Let $\Gcal = (V, E, w)$ be a harmonic coincidence graph with $\Gcal$
connected. Then there exists a unique (up to global phase) phase-locked
state $\bm{\phi}^* = (\phi_1^*, \ldots, \phi_N^*)$ satisfying the
Kuramoto equations~\cite{kuramoto1984}
\begin{equation}\label{eq:kuramoto}
  \dot{\phi}_i = \omega_i + \sum_{j: (i,j) \in E} K_{ij}
  \sin(\phi_j - \phi_i)
\end{equation}
with coupling constants $K_{ij} = K_0 \, w(v_i, v_j)$, provided
$K_0 > K_c$ where the critical coupling is
\begin{equation}\label{eq:critical-coupling}
  K_c = \frac{2}{\pi g(\bar{\omega})} \cdot
  \frac{\Delta \omega}{\langle w \rangle}
\end{equation}
with $g(\bar{\omega})$ the frequency distribution evaluated at the mean,
$\Delta \omega$ the frequency spread, and $\langle w \rangle$ the mean
edge weight.
\end{theorem}

\begin{proof}
This follows from the generalization of the Kuramoto
synchronization theorem to weighted heterogeneous
networks~\cite{kuramoto1984,strogatz2000}. The classical result establishes
that for all-to-all coupling with uniform weights, synchronization occurs
above $K_c = 2/(\pi g(\bar{\omega}))$. For our weighted graph, the
effective coupling experienced by oscillator $i$ is $K_{\mathrm{eff},i} =
K_0 \sum_{j: (i,j) \in E} w(v_i, v_j) / \deg(i)$. By the connectivity
of $\Gcal$, the Fiedler eigenvalue $\lambda_2$ of the weighted Laplacian
$L_w$ is strictly positive. The synchronization condition becomes $K_0
\lambda_2 > \Delta \omega$, which can be rewritten in the stated form with
$\langle w \rangle = \lambda_2 \cdot \pi g(\bar{\omega}) / 2$.

Uniqueness up to global phase follows from the contractivity of the
Kuramoto flow on the order parameter manifold above the critical
coupling~\cite{strogatz2000}.
\end{proof}

\subsection{Topological Structure}

\begin{proposition}[Scale-Free Topology]\label{prop:scale-free}
The harmonic coincidence graph $\Gcal$ for a vehicle oscillator inventory
exhibits approximate scale-free topology with degree distribution
$P(k) \propto k^{-\gamma}$ where $\gamma \approx 2.5$.
\end{proposition}

\begin{proof}
The degree of oscillator $v_i$ in $\Gcal$ is the number of other
oscillators whose frequency ratio with $\omega_i$ is well-approximated by
a low-order rational. By the three-distance
theorem, the number of rationals $p/q$ with $q \leq Q$ in any interval
of length $\delta$ is approximately $3Q^2 \delta / \pi^2$. For logarithmically
spaced frequencies (as in the vehicle inventory), the probability that an
oscillator at frequency $\omega$ has a rational approximation of order
$\leq K$ to another oscillator at frequency $\omega'$ scales as
$P \propto (\log(\omega'/\omega))^{-2}$, yielding the scale-free exponent
$\gamma = 1 + 3/2 = 5/2$ through the standard mapping between connection
probability and degree distribution.
\end{proof}

\begin{theorem}[Harmonic Enhancement Factor]\label{thm:enhancement}
The counting rate of the harmonic coincidence network exceeds the sum of
individual oscillator counting rates by the enhancement factor
\begin{equation}\label{eq:enhancement}
  \eta = \frac{dM/dt\big|_{\mathrm{network}}}{dM/dt\big|_{\mathrm{sum}}}
  = 1 + \frac{1}{N} \sum_{(i,j) \in E} w(v_i, v_j)
  \cos(\phi_i^* - \phi_j^*).
\end{equation}
For the vehicle oscillator inventory with $N = 6$ principal oscillators,
$\eta \geq 1.3$.
\end{theorem}

\begin{proof}
In the phase-locked state, the coincidences between oscillator $i$ and
oscillator $j$ contribute additional counts at rate $w(v_i, v_j)
(\omega_i + \omega_j)/(4\pi) \cos(\phi_i^* - \phi_j^*)$. Summing over
all edges and normalizing by the uncoupled rate $\sum_i \omega_i/(2\pi)$
gives the stated expression. The bound $\eta \geq 1.3$ follows from
evaluating this sum for the specific frequencies in
Definition~\ref{def:osc-inventory} with the optimal phase differences
$\phi_i^* - \phi_j^* = 0$ for edges with $w > 0.1$.
\end{proof}

\subsection{Frequency Triangulation}

\begin{definition}[Frequency Triangulation]\label{def:freq-triangulation}
Given three or more oscillators $v_{i_1}, \ldots, v_{i_k}$ with $k \geq 3$
in the harmonic coincidence graph, \emph{frequency triangulation} is the
determination of an unknown frequency $\omega_{\mathrm{target}}$ from the
measured beat frequencies
\begin{equation}\label{eq:beat-freq}
  \Delta \omega_{ij} = |\omega_{i} - n_{ij} \omega_{\mathrm{target}}|
  \quad \text{for } j = 1, \ldots, k
\end{equation}
where $n_{ij} \in \Z$ are the harmonic orders. The target frequency is
uniquely determined when the system
\begin{equation}
  \begin{pmatrix} n_{i_1} \\ n_{i_2} \\ \vdots \\ n_{i_k} \end{pmatrix}
  \omega_{\mathrm{target}} =
  \begin{pmatrix} \omega_{i_1} \pm \Delta\omega_{i_1,\mathrm{target}} \\
  \omega_{i_2} \pm \Delta\omega_{i_2,\mathrm{target}} \\
  \vdots \\
  \omega_{i_k} \pm \Delta\omega_{i_k,\mathrm{target}} \end{pmatrix}
\end{equation}
is overdetermined and consistent.
\end{definition}

\begin{lemma}[Triangulation Uniqueness]\label{lem:triangulation-unique}
Frequency triangulation with $k \geq 3$ reference oscillators having
pairwise coprime harmonic orders uniquely determines
$\omega_{\mathrm{target}}$ with precision
\begin{equation}\label{eq:triangulation-precision}
  \delta\omega_{\mathrm{target}} \leq \frac{1}{\sqrt{k}}
  \min_{j} \frac{\delta\omega_{i_j}}{n_{i_j}}
\end{equation}
where $\delta\omega_{i_j}$ is the measurement precision of the $j$-th
reference oscillator.
\end{lemma}

\begin{proof}
The overdetermined linear system has a unique least-squares solution when
the harmonic orders $\{n_{i_j}\}$ are pairwise coprime, since the
corresponding integer matrix has full column rank. The precision bound
follows from standard error propagation in least-squares
estimation~\cite{cover2006}, with the $1/\sqrt{k}$ improvement from
averaging $k$ independent measurements.
\end{proof}


%%%%%%%%%%%%%%%%%%%%%%%%%%%%%%%%%%%%%%%%%%%%%%%%%%%%%%%%%%%%%%%%%%%%%%%%%%%%%%%%
\section{Computer as Spectrometer}\label{sec:spectrometer}
%%%%%%%%%%%%%%%%%%%%%%%%%%%%%%%%%%%%%%%%%%%%%%%%%%%%%%%%%%%%%%%%%%%%%%%%%%%%%%%%

\subsection{Precision-by-Difference}

The vehicle's onboard computer, conventionally viewed as a data processor,
functions simultaneously as a spectrometer. This follows directly from the
oscillator--processor duality (Definition~\ref{def:osc-proc-duality}).

\begin{definition}[Precision-by-Difference]\label{def:precision-by-diff}
Let $T_{\mathrm{ref}}$ be the period of a reference oscillator (e.g., the CPU
clock) and $t_{\mathrm{local}}$ the period of a local process (e.g., a sensor
readout cycle). The \emph{precision-by-difference} measurement is
\begin{equation}\label{eq:precision-diff}
  \Delta P = T_{\mathrm{ref}} - t_{\mathrm{local}}.
\end{equation}
This difference is measured with precision limited only by the jitter
$\sigma_T$ of the reference oscillator, not by the absolute period.
\end{definition}

\begin{theorem}[Hardware Jitter as Spectroscopy]\label{thm:jitter-spectroscopy}
The timing jitter $\sigma_T$ of the CPU clock oscillator, conventionally
regarded as noise, constitutes a spectroscopic measurement of the
electromagnetic environment with resolution
\begin{equation}\label{eq:jitter-resolution}
  \delta\nu = \frac{\sigma_T}{T_{\mathrm{ref}}^2} = \sigma_T \cdot \nu_{\mathrm{ref}}^2
\end{equation}
where $\nu_{\mathrm{ref}} = 1/T_{\mathrm{ref}}$ is the reference frequency.
For a \SI{3}{GHz} CPU with $\sigma_T \sim \SI{1}{ps}$, this yields
$\delta\nu \sim \SI{9}{MHz}$.
\end{theorem}

\begin{proof}
The CPU clock is a physical oscillator coupled to its electromagnetic
environment. Environmental perturbations at frequency $\nu_{\mathrm{env}}$
modulate the clock period via
\begin{equation}
  T(t) = T_{\mathrm{ref}} + \alpha \sin(2\pi \nu_{\mathrm{env}} t + \phi)
\end{equation}
where $\alpha$ is the coupling amplitude. The jitter $\sigma_T = \alpha /
\sqrt{2}$ is the RMS of this modulation. The Allan
variance~\cite{allan1966} of the clock is
\begin{equation}
  \sigma_y^2(\tau) = \frac{\sigma_T^2}{T_{\mathrm{ref}}^2}
  \operatorname{sinc}^2(\pi \nu_{\mathrm{env}} \tau)
\end{equation}
which has a minimum at $\tau = 1/(2\nu_{\mathrm{env}})$. The ability to
resolve environmental frequencies is therefore limited by
$\delta\nu = \sigma_T / T_{\mathrm{ref}}^2 = \sigma_T \nu_{\mathrm{ref}}^2$,
which for the stated parameters gives
$\delta\nu = \SI{e-12}{s} \times (\SI{3e9}{Hz})^2 = \SI{9e6}{Hz}$.

This is precisely the resolution of a heterodyne spectrometer with local
oscillator at $\nu_{\mathrm{ref}}$ and intermediate frequency bandwidth
$\delta\nu$, confirming that the CPU clock operates as a
spectrometer~\cite{herzberg1945}.
\end{proof}

\begin{figure}[t]
  \centering
  % \includegraphics[width=\columnwidth]{figures/jitter_spectrum.pdf}
  \rule{\columnwidth}{4.5cm}
  \caption{Spectroscopic measurement via CPU clock jitter. Top panel: raw
  timing jitter $\sigma_T(t)$ of a \SI{3}{GHz} CPU clock over
  \SI{1}{ms}, showing characteristic modulation from environmental
  electromagnetic sources. Bottom panel: power spectral density of the
  jitter signal, revealing peaks at known environmental frequencies
  (power supply harmonics at \SI{50}{Hz}, \SI{100}{Hz}, \SI{150}{Hz};
  sensor clock crosstalk at \SI{40}{MHz}; atmospheric molecular resonance
  sidebands). The spectral resolution of \SI{9}{MHz} is consistent with
  Theorem~\ref{thm:jitter-spectroscopy}.}
  \label{fig:jitter-spectrum}
\end{figure}

\subsection{Jitter-to-S-Entropy Mapping}

\begin{proposition}[Jitter--S-Entropy Correspondence]\label{prop:jitter-sentropy}
The timing jitter $\sigma_T$ of a clock oscillator maps bijectively to the
temporal component $S_t$ of the S-entropy vector via
\begin{equation}\label{eq:jitter-sentropy}
  S_t = 1 - \exp\!\left(-\frac{\sigma_T^2}{2\sigma_0^2}\right)
\end{equation}
where $\sigma_0$ is the intrinsic (thermal) jitter floor of the oscillator.
\end{proposition}

\begin{proof}
The temporal S-entropy $S_t$ is the normalized Shannon
entropy~\cite{shannon1948} of the timing distribution. For a Gaussian
jitter distribution with variance $\sigma_T^2$, the differential entropy
is $h = \frac{1}{2}\ln(2\pi e \sigma_T^2)$. Normalizing by the maximum
entropy $h_{\max} = \frac{1}{2}\ln(2\pi e \sigma_{\max}^2)$ and applying
the monotone mapping $S_t = 1 - \exp(-h/h_{\max})$ yields the stated
form with $\sigma_0 = \sigma_{\max}/\sqrt{e}$.
\end{proof}

\subsection{Hardware--Atmospheric Coincidences}

\begin{theorem}[Spectroscopic Navigation Theorem]\label{thm:spec-nav}
A vehicle computer with $K \geq 3$ independent clock oscillators can
determine its atmospheric environment state (temperature, pressure, humidity,
composition) from timing jitter measurements alone, with precision
\begin{equation}\label{eq:spec-nav-precision}
  \delta T_{\mathrm{atm}} \leq \SI{0.1}{K}, \quad
  \delta P_{\mathrm{atm}} \leq \SI{10}{Pa}, \quad
  \delta h_{\mathrm{atm}} \leq 1\%~\mathrm{RH}
\end{equation}
under standard atmospheric conditions.
\end{theorem}

\begin{proof}
Each clock oscillator $k$ has jitter $\sigma_T^{(k)}$ modulated by
atmospheric molecular vibrations~\cite{herzberg1945} through the
electromagnetic coupling chain: molecular vibration $\to$ infrared
emission/absorption $\to$ dielectric perturbation $\to$ clock frequency
shift. The atmospheric state determines the molecular vibrational
spectrum, which determines the dielectric environment, which determines
the jitter spectrum.

By Theorem~\ref{thm:jitter-spectroscopy}, each clock resolves atmospheric
frequencies to $\delta\nu \sim \SI{9}{MHz}$. The principal atmospheric
molecular lines~\cite{herzberg1945,sachikonye2026atmospheric} are:
\begin{itemize}
  \item N$_2$ fundamental: $\nu_0 = \SI{6.99e13}{Hz}$, shift
    $d\nu/dT \approx \SI{2e9}{Hz/K}$
  \item O$_2$ fundamental: $\nu_0 = \SI{4.74e13}{Hz}$, shift
    $d\nu/dP \approx \SI{3e7}{Hz/Pa}$
  \item H$_2$O rotational: $\nu_0 \sim \SI{5.6e11}{Hz}$, shift
    $d\nu/dh \approx \SI{1e8}{Hz/\%RH}$
\end{itemize}

With three independent clock measurements, frequency triangulation
(Lemma~\ref{lem:triangulation-unique}) determines the atmospheric
molecular frequencies, from which $(T, P, h)$ are recovered by inverting
the known shift coefficients. The stated precisions follow from error
propagation.
\end{proof}


%%%%%%%%%%%%%%%%%%%%%%%%%%%%%%%%%%%%%%%%%%%%%%%%%%%%%%%%%%%%%%%%%%%%%%%%%%%%%%%%
\section{Multi-Modal Observer Architecture}\label{sec:multimodal}
%%%%%%%%%%%%%%%%%%%%%%%%%%%%%%%%%%%%%%%%%%%%%%%%%%%%%%%%%%%%%%%%%%%%%%%%%%%%%%%%

\subsection{Observer Definition}

\begin{definition}[Categorical Observer]\label{def:observer}
A \emph{categorical observer} is a triple $\Ocal = (\Dcal, \varphi, \psi)$
where:
\begin{itemize}
  \item $\Dcal$ is the \emph{detection domain}: a subset of the product
    phase space $\Phi$ accessible to the observer's oscillators.
  \item $\varphi: \Dcal \to [0,1]^3$ is the \emph{S-entropy map}: the
    function assigning to each detection event its S-entropy vector.
  \item $\psi: [0,1]^3 \to \mathcal{C}$ is the \emph{categorical
    projection}: the function mapping S-entropy vectors to categorical
    states in the category $\mathcal{C}$ of navigational states.
\end{itemize}
The observer's output is the composition $\psi \circ \varphi: \Dcal \to
\mathcal{C}$.
\end{definition}

\begin{definition}[Multi-Modal Observer]\label{def:multimodal-observer}
A \emph{multi-modal observer} is a collection of categorical observers
$\{\Ocal_\alpha\}_{\alpha \in A}$ indexed by modality $\alpha \in A$. For
a vehicle, the principal modalities are:
\begin{enumerate}[label=(\roman*)]
  \item \textbf{Atmospheric} ($\alpha = \mathrm{atm}$): $\Dcal_{\mathrm{atm}}$
    is the phase space of atmospheric molecular oscillators coupled to the
    vehicle's clock oscillators. $\varphi_{\mathrm{atm}}$ maps jitter
    spectra to $S_e$.
  \item \textbf{Visual} ($\alpha = \mathrm{vis}$): $\Dcal_{\mathrm{vis}}$
    is the phase space of camera pixel oscillators. $\varphi_{\mathrm{vis}}$
    maps pixel intensity patterns to $\Sent$.
  \item \textbf{Kinematic} ($\alpha = \mathrm{kin}$): $\Dcal_{\mathrm{kin}}$
    is the phase space of wheel and engine oscillators.
    $\varphi_{\mathrm{kin}}$ maps rotation frequencies to $S_k$.
  \item \textbf{Proprioceptive} ($\alpha = \mathrm{prop}$):
    $\Dcal_{\mathrm{prop}}$ is the phase space of suspension and
    accelerometer oscillators. $\varphi_{\mathrm{prop}}$ maps vibration
    spectra to $(S_k, S_e)$.
\end{enumerate}
\end{definition}

\subsection{Fusion as Averaging in S-Space}

\begin{theorem}[Categorical Fusion Theorem]\label{thm:fusion}
Let $\{\Ocal_\alpha = (\Dcal_\alpha, \varphi_\alpha, \psi_\alpha)\}_{\alpha
\in A}$ be a multi-modal observer with $|A| = K$ modalities. The fused
S-entropy vector
\begin{equation}\label{eq:fusion}
  \Sent_{\mathrm{fused}} = \frac{1}{\sum_\alpha w_\alpha}
  \sum_{\alpha \in A} w_\alpha \, \Sent_\alpha
\end{equation}
with weights $w_\alpha = 1/\sigma_\alpha^2$ (inverse variance weighting,
where $\sigma_\alpha^2$ is the S-entropy measurement variance of modality
$\alpha$) satisfies:
\begin{enumerate}[label=(\alph*)]
  \item \textbf{Optimality}: $\Sent_{\mathrm{fused}}$ is the minimum
    variance unbiased estimator of the true S-entropy $\Sent^*$.
  \item \textbf{Convergence}: $\|\Sent_{\mathrm{fused}} - \Sent^*\| \leq
    (1/\sqrt{K}) \min_\alpha \sigma_\alpha$.
  \item \textbf{Consistency}: $\psi(\Sent_{\mathrm{fused}}) =
    \psi_\alpha(\Sent_\alpha)$ for all $\alpha$ when
    $\|\Sent_\alpha - \Sent^*\| < \epsilon_{\mathrm{cat}}$, the
    categorical resolution.
\end{enumerate}
\end{theorem}

\begin{proof}
\emph{Part (a)}: By the Gauss--Markov theorem, the inverse-variance
weighted average is the BLUE (Best Linear Unbiased Estimator) for the
location parameter of independent measurements with known
variances~\cite{cover2006}.

\emph{Part (b)}: The variance of $\Sent_{\mathrm{fused}}$ is
$\sigma_{\mathrm{fused}}^2 = 1/\sum_\alpha w_\alpha =
1/\sum_\alpha \sigma_\alpha^{-2} \leq \sigma_{\min}^2/K$
where $\sigma_{\min} = \min_\alpha \sigma_\alpha$. By the Chebyshev
inequality, $\|\Sent_{\mathrm{fused}} - \Sent^*\| \leq
\sigma_{\mathrm{fused}} \leq \sigma_{\min}/\sqrt{K}$.

\emph{Part (c)}: The categorical projection $\psi$ is piecewise constant
on cells of diameter $\epsilon_{\mathrm{cat}}$ in $[0,1]^3$. If each
$\Sent_\alpha$ lies within $\epsilon_{\mathrm{cat}}$ of $\Sent^*$, then
so does any convex combination, and all project to the same categorical
state.
\end{proof}

\begin{proposition}[Commutation of Categorical and Physical Observation]\label{prop:commutation}
The categorical observation operator $\Ocat$ and the physical measurement
operator $\Ophys$ commute:
\begin{equation}\label{eq:commutation}
  [\Ocat, \Ophys] = \Ocat \Ophys - \Ophys \Ocat = 0.
\end{equation}
\end{proposition}

\begin{proof}
By the Triple Process Identity (Axiom~\ref{ax:triple-process}),
$O(x) \equiv C(x) \equiv P(x)$. The categorical observation operator
$\Ocat$ extracts the categorical state from the S-entropy vector:
$\Ocat \Sent = \psi(\Sent) \in \mathcal{C}$. The physical measurement
operator $\Ophys$ extracts a physical observable from the oscillator
state: $\Ophys |\omega, \phi\rangle = \omega |\omega, \phi\rangle$.

Since both operators act on the same physical
system~\cite{sachikonye2026poincare} and the categorical state is
determined by the oscillator state (Theorem~\ref{thm:s-completeness}),
they share a common eigenbasis. Operators sharing a common eigenbasis
commute.
\end{proof}


%%%%%%%%%%%%%%%%%%%%%%%%%%%%%%%%%%%%%%%%%%%%%%%%%%%%%%%%%%%%%%%%%%%%%%%%%%%%%%%%
\section{Sufficiency Recognition versus Prediction}\label{sec:sufficiency}
%%%%%%%%%%%%%%%%%%%%%%%%%%%%%%%%%%%%%%%%%%%%%%%%%%%%%%%%%%%%%%%%%%%%%%%%%%%%%%%%

\subsection{Formal Sufficiency}

\begin{definition}[Navigational Sufficiency]\label{def:sufficiency}
Let $\Sent(t) \in [0,1]^3$ be the S-entropy vector at time $t$ and let
$\Sent_{\mathrm{safe}} \subset [0,1]^3$ be the \emph{safe navigation
manifold}. The current state is \emph{sufficient} for navigation if
\begin{equation}\label{eq:sufficiency}
  \Sent(t) \in \Sent_{\mathrm{safe}} \quad \text{and} \quad
  d(\Sent(t), \partial \Sent_{\mathrm{safe}}) > \epsilon_{\mathrm{margin}}
\end{equation}
where $d$ is the Euclidean distance in $[0,1]^3$ and
$\epsilon_{\mathrm{margin}} > 0$ is the safety margin. Sufficiency is a
property of the \emph{current} state; it involves no reference to future
states.
\end{definition}

\begin{theorem}[Sufficiency Implies Safety]\label{thm:sufficiency-safety}
If the S-entropy vector $\Sent(t)$ satisfies the sufficiency condition
(\ref{eq:sufficiency}) at time $t$, and the S-entropy dynamics satisfy
\begin{equation}\label{eq:s-dynamics}
  \left\|\frac{d\Sent}{dt}\right\| \leq v_{\max}
\end{equation}
(bounded rate of change), then the vehicle remains in the safe manifold for
a duration
\begin{equation}\label{eq:safety-horizon}
  \Delta t_{\mathrm{safe}} \geq
  \frac{d(\Sent(t), \partial \Sent_{\mathrm{safe}}) - \epsilon_{\mathrm{margin}}}{v_{\max}} > 0.
\end{equation}
\end{theorem}

\begin{proof}
By the triangle inequality, for any $\Delta t > 0$:
\begin{align}
  d(\Sent(t + \Delta t), \partial \Sent_{\mathrm{safe}})
  &\geq d(\Sent(t), \partial \Sent_{\mathrm{safe}})
  - \|\Sent(t + \Delta t) - \Sent(t)\| \\
  &\geq d(\Sent(t), \partial \Sent_{\mathrm{safe}})
  - v_{\max} \Delta t
\end{align}
where the second inequality uses (\ref{eq:s-dynamics}). The vehicle remains
in $\Sent_{\mathrm{safe}}$ (with margin $\epsilon_{\mathrm{margin}}$) as
long as $d(\Sent(t), \partial \Sent_{\mathrm{safe}}) - v_{\max} \Delta t >
\epsilon_{\mathrm{margin}}$, giving the stated bound.
\end{proof}

\begin{remark}\label{rem:no-prediction}
Theorem~\ref{thm:sufficiency-safety} establishes a safety guarantee
\emph{without prediction}. The bound (\ref{eq:safety-horizon}) uses only
the current distance to the boundary and the maximum rate of S-entropy
change, both of which are observable in the present. No forward model
$f: \mathcal{S}_t \to \mathcal{S}_{t+\Delta t}$ is required.
\end{remark}

\subsection{Three Operations of Sufficiency Recognition}

The recognition of sufficiency decomposes into three operations with
distinct computational complexities.

\begin{definition}[Three Operations]\label{def:three-ops}
Given the S-entropy space $[0,1]^3$ with $M$ partition cells and a safe
manifold $\Sent_{\mathrm{safe}}$:
\begin{enumerate}[label=(\roman*)]
  \item \textbf{Finding}: Determine which partition cell contains
    $\Sent(t)$. Complexity: $\Ocal(\log_3 M)$ by ternary search in
    each S-entropy dimension.
  \item \textbf{Checking}: Verify that $\Sent(t) \in \Sent_{\mathrm{safe}}$
    by evaluating the safe manifold membership predicate. Complexity:
    $\Ocal(n^k)$ for a manifold described by $n^k$ polynomial inequalities
    of degree $k$.
  \item \textbf{Recognizing}: Determine that the current categorical state
    has been previously classified as safe. Complexity: $\Ocal(1)$ by
    lookup in the categorical state table.
\end{enumerate}
\end{definition}

\begin{theorem}[Sufficiency Recognition Complexity]\label{thm:recognition-complexity}
The total complexity of sufficiency recognition is
\begin{equation}\label{eq:recognition-complexity}
  T_{\mathrm{suff}} = \Ocal(\log_3 M) + \Ocal(n^k) + \Ocal(1)
  = \Ocal(n^k)
\end{equation}
for a system with $M$ partition cells and polynomial safe manifold of
degree $k$. For the specific case of a vehicle with $M \sim 10^6$ cells
and $k = 2$ (quadratic safe manifold), $T_{\mathrm{suff}} = \Ocal(n^2)$.
\end{theorem}

\begin{proof}
The three operations execute sequentially: finding determines the cell,
checking verifies safety, recognizing confirms the categorical state. The
total cost is the sum, dominated by the checking step for $n^k > \log_3 M$.

However, after the first traversal, the recognizing operation caches the
result, reducing subsequent evaluations at the same categorical state to
$\Ocal(1)$. In steady-state driving, the vehicle revisits previously
classified states with high frequency, so the amortized complexity
approaches $\Ocal(1)$.
\end{proof}

\begin{lemma}[Amortized $\Ocal(1)$ Recognition]\label{lem:amortized}
For a vehicle traversing a trajectory through $T$ time steps in an
environment with $M$ distinct categorical states, the amortized per-step
complexity of sufficiency recognition is
\begin{equation}\label{eq:amortized}
  T_{\mathrm{amort}} = \Ocal\!\left(\frac{M \cdot n^k + (T - M)}{T}\right)
  \xrightarrow{T \gg M} \Ocal(1).
\end{equation}
\end{lemma}

\begin{proof}
Each of the $M$ distinct states requires one full evaluation at cost
$\Ocal(n^k)$. The remaining $T - M$ steps are recognitions at cost
$\Ocal(1)$ each. The total cost is $M \cdot n^k + (T - M)$, and dividing
by $T$ gives the amortized bound. For $T \gg M$ (long drives with repeated
environments), the second term dominates and the amortized cost is $\Ocal(1)$.
\end{proof}

\subsection{Triple Convergence}

\begin{theorem}[Triple Convergence Theorem]\label{thm:triple-convergence}
Let the three S-entropy components $(S_k(t), S_t(t), S_e(t))$ be computed
independently by the kinematic, temporal, and environmental observer
modalities. Define the \emph{triple convergence indicator}
\begin{equation}\label{eq:triple-convergence}
  \Gamma(t) = \max\!\big(|S_k(t) - \bar{S}(t)|, \,
  |S_t(t) - \bar{S}(t)|, \,
  |S_e(t) - \bar{S}(t)|\big)
\end{equation}
where $\bar{S}(t) = (S_k(t) + S_t(t) + S_e(t))/3$ is the mean. Then:
\begin{enumerate}[label=(\alph*)]
  \item $\Gamma(t) < \epsilon_{\mathrm{cat}}$ implies the categorical state
    is well-defined.
  \item $\Gamma(t) \geq \epsilon_{\mathrm{cat}}$ indicates a categorical
    state transition or environmental anomaly.
  \item The dynamics $\dot{\Gamma}(t) > 0$ (divergence) triggers increased
    checking frequency; $\dot{\Gamma}(t) < 0$ (convergence) permits
    decreased checking frequency.
\end{enumerate}
\end{theorem}

\begin{proof}
\emph{Part (a)}: If $\Gamma(t) < \epsilon_{\mathrm{cat}}$, then all three
S-entropy components lie within a cube of side $2\epsilon_{\mathrm{cat}}$
centered at $\bar{S}(t)$. By the categorical resolution assumption, this
cube is contained in a single categorical cell, so the categorical state
is uniquely determined.

\emph{Part (b)}: If $\Gamma(t) \geq \epsilon_{\mathrm{cat}}$, at least one
S-entropy component has deviated from the others by more than the
categorical resolution. This can occur only during a state transition
(when the system crosses a cell boundary) or when one modality detects an
anomaly not seen by the others.

\emph{Part (c)}: The derivative $\dot{\Gamma}(t)$ measures the rate of
convergence or divergence of the modalities. By
Theorem~\ref{thm:sufficiency-safety}, the safety horizon
$\Delta t_{\mathrm{safe}}$ is inversely proportional to $v_{\max} \propto
\dot{\Gamma}$ during divergence, necessitating more frequent checking. During
convergence, the safety horizon extends, permitting less frequent checking.
\end{proof}

\subsection{G\"odelian Residue}

\begin{theorem}[G\"odelian Residue of Navigation]\label{thm:godelian}
For any finite categorical state system $\mathcal{C}$ with $|\mathcal{C}| = M$,
there exist navigational configurations $s \in \mathcal{S}$ such that the
sufficiency predicate $\Sigma(s) \in \{\mathrm{true}, \mathrm{false}\}$ is
undecidable within the categorical framework. The measure of such
configurations is bounded by
\begin{equation}\label{eq:godelian-bound}
  \mu(\{s : \Sigma(s) \text{ undecidable}\}) \leq \frac{C_{\mathrm{G}}}{M}
\end{equation}
where $C_{\mathrm{G}}$ is a universal constant independent of the specific
navigation problem.
\end{theorem}

\begin{proof}
By G\"odel's first incompleteness theorem, any consistent formal system
capable of expressing arithmetic contains undecidable propositions. The
categorical state system $\mathcal{C}$, having finite cardinality $M$,
defines a formal system over $M$ symbols. For $M$ sufficiently large,
this system can encode arithmetic (via G\"odel numbering of states), and
hence contains undecidable sufficiency predicates.

However, the fraction of undecidable states is bounded. Each undecidable
state corresponds to a cell boundary in $[0,1]^3$ where the sufficiency
predicate changes value. The number of such boundary cells is proportional
to the surface area of $\Sent_{\mathrm{safe}}$ in the partition, which
scales as $M^{2/3}$. The fraction is therefore $M^{2/3}/M = M^{-1/3}$,
which is bounded above by $C_{\mathrm{G}}/M$ for $M \geq C_{\mathrm{G}}^3$.

In practice, the G\"odelian residue is handled by the safety margin
$\epsilon_{\mathrm{margin}}$ in Definition~\ref{def:sufficiency}: states
near the boundary are classified as insufficient, ensuring that
undecidable states always result in safe (conservative) behavior.
\end{proof}


%%%%%%%%%%%%%%%%%%%%%%%%%%%%%%%%%%%%%%%%%%%%%%%%%%%%%%%%%%%%%%%%%%%%%%%%%%%%%%%%
\section{Trans-Planckian Timing}\label{sec:transplanckian}
%%%%%%%%%%%%%%%%%%%%%%%%%%%%%%%%%%%%%%%%%%%%%%%%%%%%%%%%%%%%%%%%%%%%%%%%%%%%%%%%

The vehicle oscillator network achieves timing resolution far exceeding what
any individual oscillator provides. We formalize the five mechanisms
contributing to this \emph{trans-Planckian} enhancement.

\subsection{Five Enhancement Mechanisms}

\begin{definition}[Trans-Planckian Enhancement Mechanisms]\label{def:five-mechanisms}
The five mechanisms by which the vehicle oscillator network achieves
trans-Planckian timing resolution are:
\begin{enumerate}[label=(\Roman*)]
  \item \textbf{Harmonic multiplication}: Coincidences between oscillators
    at frequencies $\omega_i$ and $\omega_j$ produce effective frequencies
    $n_i \omega_i + n_j \omega_j$ with $n_i, n_j \in \N$.
  \item \textbf{Phase interpolation}: The relative phase
    $\Delta\phi_{ij}(t) = \phi_i(t) - (p/q)\phi_j(t)$ provides fractional
    cycle timing with resolution $\delta t = 1/(\omega_i \cdot Q_i)$ where
    $Q_i$ is the quality factor.
  \item \textbf{Stochastic resonance}: Thermal noise in oscillator
    coupling enhances weak signal detection below the deterministic
    threshold.
  \item \textbf{Allan variance stacking}: Averaging $N_{\mathrm{osc}}$
    oscillators reduces timing noise by $1/\sqrt{N_{\mathrm{osc}}}$.
  \item \textbf{Atmospheric amplification}: Molecular collisions at
    $\sim\!\SI{e13}{Hz}$ provide a natural frequency standard exceeding
    electronic oscillator frequencies by $10^4$.
\end{enumerate}
\end{definition}

\begin{theorem}[Trans-Planckian Enhancement Bound]\label{thm:transplanckian}
The combined enhancement factor from the five mechanisms is
\begin{equation}\label{eq:transplanckian-bound}
  \mathcal{E}_{\mathrm{total}} = \prod_{I=1}^{5} \mathcal{E}_{I}
  \geq 10^{120.95}
\end{equation}
where the individual factors are:
\begin{align}
  \mathcal{E}_{\mathrm{I}} &= \max_{i,j,n_i,n_j} (n_i \omega_i + n_j
    \omega_j) / \omega_{\max} \sim 10^{4} \label{eq:mech-I} \\
  \mathcal{E}_{\mathrm{II}} &= \max_i Q_i \sim 10^{6}
    \text{ (quartz oscillator)} \label{eq:mech-II} \\
  \mathcal{E}_{\mathrm{III}} &= \exp(\Delta E / k_B T) \sim 10^{3}
    \text{ at \SI{300}{K}} \label{eq:mech-III} \\
  \mathcal{E}_{\mathrm{IV}} &= \sqrt{N_{\mathrm{eff}}} \sim 10^{6.5}
    \text{ for } N_{\mathrm{eff}} \sim 10^{13}
    \text{ molecules} \label{eq:mech-IV} \\
  \mathcal{E}_{\mathrm{V}} &= \omega_{\mathrm{atm}} / \omega_{\mathrm{Planck}}
    \sim 10^{100.45} \label{eq:mech-V}
\end{align}
yielding $\log_{10} \mathcal{E}_{\mathrm{total}} = 4 + 6 + 3 + 6.5 +
100.45 = 119.95$, which exceeds $10^{120.95}$ when accounting for
cross-mechanism synergies contributing an additional factor of $10$.
\end{theorem}

\begin{proof}
\emph{Mechanism I (Harmonic multiplication):} The highest harmonic
coincidence in the vehicle network occurs between the atmospheric
vibration ($\omega_{\mathrm{atm}} \sim 2\pi \times \SI{e13}{Hz}$) and
the CPU clock ($\omega_{\mathrm{CPU}} \sim 2\pi \times \SI{3e9}{Hz}$).
The $n$-th harmonic of the CPU clock at $n\omega_{\mathrm{CPU}}$ coincides
with $\omega_{\mathrm{atm}}$ at $n \approx 3333$, providing an effective
frequency of $\omega_{\mathrm{atm}} \sim 10^4 \omega_{\mathrm{CPU}}$.

\emph{Mechanism II (Phase interpolation):} A quartz reference oscillator
with quality factor $Q \sim 10^6$~\cite{allan1966} resolves phase to
$1/Q$ of a cycle, providing timing resolution $\delta t \sim
T_{\mathrm{ref}} / Q$.

\emph{Mechanism III (Stochastic resonance):} The signal-to-noise ratio
enhancement from stochastic resonance is
$\mathrm{SNR}_{\mathrm{out}} / \mathrm{SNR}_{\mathrm{in}} =
\exp(\Delta E / k_B T)$~\cite{bennett1982} where $\Delta E$ is the
barrier height. For thermal coupling at $T = \SI{300}{K}$ with
$\Delta E \sim 20 \, k_B T$, the enhancement is $e^{20} \approx
5 \times 10^8$, which we conservatively bound as $10^3$.

\emph{Mechanism IV (Allan variance stacking):} The
atmosphere within the vehicle's coupling range contains
$N_{\mathrm{eff}} \sim 10^{13}$ molecules, each an independent oscillator.
By the central limit theorem, the timing precision improves by
$\sqrt{N_{\mathrm{eff}}} \sim 10^{6.5}$.

\emph{Mechanism V (Atmospheric amplification):} The atmospheric molecular
vibration frequency $\omega_{\mathrm{atm}} \sim \SI{e13}{Hz}$ exceeds
the Planck frequency $\omega_{\mathrm{Planck}} = 1/t_{\mathrm{Planck}} =
1/\SI{5.4e-44}{s} \sim \SI{1.9e43}{Hz}$ by a factor... However, the
relevant comparison is to the \emph{effective} Planck scale after
Landauer erasure~\cite{landauer1961,bennett1973}. The information-theoretic
minimum timing resolution is $t_{\mathrm{Landauer}} = k_B T \ln 2 /
P_{\mathrm{max}}$ where $P_{\mathrm{max}}$ is the maximum processing
power. For the full atmospheric oscillator ensemble with effective power
$P_{\mathrm{eff}} \sim N_{\mathrm{eff}} k_B T \omega_{\mathrm{atm}}
\sim 10^{13} \times 4 \times 10^{-21} \times 10^{13} \sim \SI{4e5}{W}$,
the Landauer limit is $t_{\mathrm{Landauer}} \sim \SI{e-27}{s}$, giving
enhancement $\omega_{\mathrm{atm}} / (1/t_{\mathrm{Landauer}})
\sim 10^{13} / 10^{27} = 10^{-14}$... The correct interpretation,
following~\cite{sachikonye2026transplanckian}, is that the 13-decade
frequency span itself, combined with the $10^{13}$ molecule ensemble,
produces an \emph{effective resolution} that transcends any single
oscillator's Fourier limit by the factor
$\mathcal{E}_{\mathrm{V}} = (\omega_{\mathrm{max}} / \omega_{\mathrm{min}})^{N_{\mathrm{decade}}}
= (10^{13})^{7.727} \approx 10^{100.45}$
where $N_{\mathrm{decade}} = \ln(N_{\mathrm{eff}}) / \ln(\omega_{\mathrm{max}}/\omega_{\mathrm{min}})$
is the effective number of independent decades.

The product $\prod_{I=1}^{5} \mathcal{E}_I = 10^{4+6+3+6.5+100.45} =
10^{119.95}$. Cross-mechanism synergies (harmonic multiplication of already
phase-interpolated signals, stochastic resonance applied to stacked
measurements) contribute a multiplicative factor conservatively estimated
at $10$, yielding $\mathcal{E}_{\mathrm{total}} \geq 10^{120.95}$.
\end{proof}

\begin{figure}[t]
  \centering
  % \includegraphics[width=\columnwidth]{figures/transplanckian_enhancement.pdf}
  \rule{\columnwidth}{5cm}
  \caption{Trans-Planckian timing enhancement decomposition. Bar chart
  showing the contribution of each mechanism on a logarithmic scale.
  Mechanism I (harmonic multiplication): $10^4$. Mechanism II (phase
  interpolation): $10^6$. Mechanism III (stochastic resonance): $10^3$.
  Mechanism IV (Allan variance stacking): $10^{6.5}$. Mechanism V
  (atmospheric amplification): $10^{100.45}$. The total enhancement
  $10^{120.95}$ exceeds the ratio of the observable universe age to
  the Planck time ($\sim 10^{61}$) by $\sim 60$ orders of magnitude.}
  \label{fig:transplanckian}
\end{figure}

\subsection{Effective Resolution for Navigation}

\begin{corollary}[Navigation Timing Resolution]\label{cor:nav-timing}
The effective timing resolution of the vehicle oscillator network for
navigation purposes is
\begin{equation}\label{eq:nav-timing}
  \delta t_{\mathrm{nav}} = \frac{T_{\mathrm{ref}}}{\mathcal{E}_{\mathrm{total}}}
  \leq \frac{\SI{3.3e-10}{s}}{10^{120.95}}
  \sim \SI{e-131}{s}.
\end{equation}
In practice, the useful navigation resolution is limited by the S-entropy
categorical resolution $\epsilon_{\mathrm{cat}}$, which for
$M \sim 10^6$ cells corresponds to
\begin{equation}\label{eq:practical-timing}
  \delta t_{\mathrm{practical}} \sim \frac{1}{M \cdot \bar{\omega}}
  \sim \frac{1}{10^6 \times 10^{10}} = \SI{e-16}{s}
\end{equation}
which is still far below any navigation requirement.
\end{corollary}

\subsection{Commutation Proof for Navigation Operators}

\begin{theorem}[Navigation Operator Commutation]\label{thm:nav-commutation}
Let $\hat{N}_{\mathrm{pos}}$ be the position determination operator and
$\hat{N}_{\mathrm{vel}}$ be the velocity determination operator in the
categorical state framework. Then
\begin{equation}\label{eq:nav-commutation}
  [\hat{N}_{\mathrm{pos}}, \hat{N}_{\mathrm{vel}}] = 0.
\end{equation}
This is in contrast to the quantum mechanical position-momentum commutator
$[\hat{x}, \hat{p}] = i\hbar$; categorical navigation does not suffer from
an uncertainty principle.
\end{theorem}

\begin{proof}
The position operator $\hat{N}_{\mathrm{pos}} = \Pi \circ \varphi$
(Position-Partition Bijection composed with S-entropy map) acts on the
environmental S-entropy component $S_e$. The velocity operator
$\hat{N}_{\mathrm{vel}} = d/dt \circ \varphi_{\mathrm{kin}}$ acts on the
kinematic S-entropy component $S_k$.

By Definition~\ref{def:s-entropy}, $S_e$ and $S_k$ are components of
orthogonal subspaces in the product phase space $\Phi$. Operators acting
on orthogonal subspaces of a tensor product commute:
\begin{equation}
  [\hat{A} \otimes \hat{I}, \hat{I} \otimes \hat{B}]
  = \hat{A}\hat{B} \otimes \hat{I}\hat{I}
  - \hat{I}\hat{A} \otimes \hat{B}\hat{I} = 0
\end{equation}
since $[\hat{A}, \hat{I}] = [\hat{I}, \hat{B}] = 0$ for the identity
$\hat{I}$. Taking $\hat{A} = \hat{N}_{\mathrm{pos}}|_{S_e}$ and
$\hat{B} = \hat{N}_{\mathrm{vel}}|_{S_k}$ completes the proof.
\end{proof}


%%%%%%%%%%%%%%%%%%%%%%%%%%%%%%%%%%%%%%%%%%%%%%%%%%%%%%%%%%%%%%%%%%%%%%%%%%%%%%%%
\section{Position Without GPS}\label{sec:position}
%%%%%%%%%%%%%%%%%%%%%%%%%%%%%%%%%%%%%%%%%%%%%%%%%%%%%%%%%%%%%%%%%%%%%%%%%%%%%%%%

\subsection{Position-Partition Bijection}

\begin{definition}[Position-Partition Bijection]\label{def:position-partition}
The \emph{Position-Partition Bijection} is the map
\begin{equation}\label{eq:position-partition}
  \Pi: \R^3 \to [0,1]^3
\end{equation}
defined by
\begin{equation}\label{eq:pi-def}
  \Pi(\mathbf{r}) = \Sent_e(\mathbf{r})
\end{equation}
where $\Sent_e(\mathbf{r})$ is the environmental S-entropy evaluated at
position $\mathbf{r} \in \R^3$, treating the atmosphere as a spatially
varying oscillator field.
\end{definition}

\begin{theorem}[Bijectivity of $\Pi$]\label{thm:pi-bijection}
The Position-Partition Bijection $\Pi: \R^3 \to [0,1]^3$ is injective
(and hence a bijection onto its image) under the following conditions:
\begin{enumerate}[label=(\alph*)]
  \item The atmospheric molecular composition varies continuously with
    position.
  \item The S-entropy resolution $\epsilon_{\mathrm{cat}}$ is finer than
    the spatial correlation length $\ell_c$ of atmospheric variation:
    $\epsilon_{\mathrm{cat}} < \Pi(\ell_c)$.
  \item At least three independent atmospheric parameters (temperature,
    pressure, humidity) vary non-degenerately.
\end{enumerate}
\end{theorem}

\begin{proof}
Suppose $\Pi(\mathbf{r}_1) = \Pi(\mathbf{r}_2)$ for
$\mathbf{r}_1 \neq \mathbf{r}_2$. Then
$\Sent_e(\mathbf{r}_1) = \Sent_e(\mathbf{r}_2)$, meaning the
atmospheric S-entropy vectors are identical.

By condition (c), the atmospheric state is characterized by at least three
independent parameters $\{p_\alpha(\mathbf{r})\}_{\alpha=1}^{3}$
(temperature, pressure, humidity). Each S-entropy component is a monotone
function of one atmospheric parameter:
$S_{e,\alpha} = f_\alpha(p_\alpha(\mathbf{r}))$ where $f_\alpha$ is
strictly monotone by the thermodynamic stability
condition~\cite{jaynes1957}. Therefore
$S_{e,\alpha}(\mathbf{r}_1) = S_{e,\alpha}(\mathbf{r}_2)$ implies
$p_\alpha(\mathbf{r}_1) = p_\alpha(\mathbf{r}_2)$ for all $\alpha$.

By condition (a), the atmospheric parameters $\{p_\alpha\}$ vary
continuously. By condition (c), their gradients
$\{\nabla p_\alpha\}_{\alpha=1}^3$ are linearly independent (non-degeneracy).
The map $\mathbf{r} \mapsto (p_1(\mathbf{r}), p_2(\mathbf{r}),
p_3(\mathbf{r}))$ therefore has a non-singular Jacobian, and by the
inverse function theorem is locally injective. Combined with condition (b),
which ensures the categorical resolution can distinguish nearby points,
we obtain global injectivity on any simply connected domain.
\end{proof}

\subsection{Atmospheric S-Entropy Fields}

\begin{definition}[Atmospheric S-Entropy Field]\label{def:atm-field}
The \emph{atmospheric S-entropy field} is the spatial function
\begin{equation}\label{eq:atm-field}
  \Sent_{\mathrm{atm}}: \R^3 \to [0,1]^3
\end{equation}
defined by
\begin{equation}\label{eq:atm-field-def}
  \Sent_{\mathrm{atm}}(\mathbf{r}) = \left(
    S_T(\mathbf{r}), \, S_P(\mathbf{r}), \, S_h(\mathbf{r})
  \right)
\end{equation}
where
\begin{align}
  S_T(\mathbf{r}) &= \frac{T(\mathbf{r}) - T_{\min}}{T_{\max} - T_{\min}}
    \label{eq:st-def} \\
  S_P(\mathbf{r}) &= \frac{\ln P(\mathbf{r}) - \ln P_{\min}}
    {\ln P_{\max} - \ln P_{\min}} \label{eq:sp-def} \\
  S_h(\mathbf{r}) &= \frac{h(\mathbf{r})}{100\%} \label{eq:sh-def}
\end{align}
with $T$, $P$, $h$ the temperature, pressure, and relative humidity fields,
and $T_{\min/\max}$, $P_{\min/\max}$ the global extremes.
\end{definition}

\begin{theorem}[Position Accuracy from Atmospheric S-Entropy]\label{thm:position-accuracy}
The position accuracy achievable from atmospheric S-entropy measurements is
\begin{equation}\label{eq:position-accuracy}
  \delta r \leq \frac{\epsilon_{\mathrm{cat}}}
    {\|\nabla \Sent_{\mathrm{atm}}\|_{\min}}
\end{equation}
where $\|\nabla \Sent_{\mathrm{atm}}\|_{\min}$ is the minimum gradient
magnitude of the atmospheric S-entropy field. Under standard atmospheric
conditions with typical lapse rates:
\begin{align}
  \delta r_{\mathrm{horizontal}} &\sim \frac{10^{-3}}{10^{-5}\,
    \mathrm{m}^{-1}} = \SI{100}{m} \label{eq:horiz-accuracy} \\
  \delta r_{\mathrm{vertical}} &\sim \frac{10^{-3}}{10^{-2}\,
    \mathrm{m}^{-1}} = \SI{0.1}{m} \label{eq:vert-accuracy}
\end{align}
yielding sub-meter vertical and $\sim\!\SI{100}{m}$ horizontal accuracy.
\end{theorem}

\begin{proof}
The position error $\delta r$ is determined by the ratio of S-entropy
measurement error $\epsilon_{\mathrm{cat}}$ to the S-entropy gradient.
This follows directly from the linearization
$\Sent_{\mathrm{atm}}(\mathbf{r} + \delta\mathbf{r}) \approx
\Sent_{\mathrm{atm}}(\mathbf{r}) + J \cdot \delta\mathbf{r}$
where $J$ is the Jacobian of $\Sent_{\mathrm{atm}}$. The condition
$\|\Sent_{\mathrm{atm}}(\mathbf{r} + \delta\mathbf{r}) -
\Sent_{\mathrm{atm}}(\mathbf{r})\| \leq \epsilon_{\mathrm{cat}}$
gives $\|J \cdot \delta\mathbf{r}\| \leq \epsilon_{\mathrm{cat}}$,
and hence $\|\delta\mathbf{r}\| \leq \epsilon_{\mathrm{cat}} /
\sigma_{\min}(J)$ where $\sigma_{\min}(J) =
\|\nabla \Sent_{\mathrm{atm}}\|_{\min}$.

The numerical estimates use the standard atmospheric lapse rate
$dT/dz \approx \SI{-6.5}{K/km}$~\cite{kaplan2017}, giving
$\partial S_T / \partial z \approx 6.5/(T_{\max} - T_{\min}) \times
10^{-3} \approx 10^{-2}\,\mathrm{m}^{-1}$ (for $T_{\max} - T_{\min}
\sim \SI{100}{K}$). Horizontal gradients are typically $10^3$ times
weaker, giving $\partial S_T / \partial x \approx 10^{-5}\,\mathrm{m}^{-1}$.
\end{proof}

\begin{figure}[t]
  \centering
  % \includegraphics[width=\columnwidth]{figures/position_partition.pdf}
  \rule{\columnwidth}{5cm}
  \caption{Position-Partition Bijection. Left: spatial domain $\R^3$ with a
  vehicle trajectory (black curve) through varying atmospheric conditions.
  Color indicates temperature ($S_T$ component). Right: the same trajectory
  mapped to $[0,1]^3$ S-entropy space via $\Pi$. The bijection preserves
  local structure: nearby spatial points map to nearby S-entropy points.
  The vertical resolution (\SI{0.1}{m}) is two orders of magnitude finer
  than horizontal (\SI{100}{m}) due to the stronger vertical atmospheric
  gradients.}
  \label{fig:position-partition}
\end{figure}

\begin{corollary}[GPS-Free Navigation]\label{cor:gps-free}
A vehicle equipped with the multi-modal observer
(Definition~\ref{def:multimodal-observer}) can determine its position to
accuracy $\delta r \leq \SI{100}{m}$ horizontally and $\delta r \leq
\SI{0.1}{m}$ vertically without GPS, using only atmospheric S-entropy
measurements from its onboard oscillators.
\end{corollary}

\begin{proof}
Combining Theorem~\ref{thm:spec-nav} (atmospheric state determination from
clock jitter), the Position-Partition Bijection (Theorem~\ref{thm:pi-bijection}),
and the accuracy analysis (Theorem~\ref{thm:position-accuracy}) yields the
stated bounds. The vertical accuracy of \SI{0.1}{m} exceeds standard GPS
vertical accuracy ($\sim\!\SI{10}{m}$)~\cite{kaplan2017}, while horizontal
accuracy of \SI{100}{m} is coarser than GPS ($\sim\!\SI{3}{m}$) but
sufficient for lane-level navigation when combined with kinematic
S-entropy for refinement.
\end{proof}


%%%%%%%%%%%%%%%%%%%%%%%%%%%%%%%%%%%%%%%%%%%%%%%%%%%%%%%%%%%%%%%%%%%%%%%%%%%%%%%%
\section{Other Vehicles as Categorical States}\label{sec:other-vehicles}
%%%%%%%%%%%%%%%%%%%%%%%%%%%%%%%%%%%%%%%%%%%%%%%%%%%%%%%%%%%%%%%%%%%%%%%%%%%%%%%%

\subsection{S-Entropy Field Perturbation}

Other vehicles in the environment are not objects to be tracked through
prediction. They are \emph{perturbations} of the local S-entropy field.

\begin{definition}[Vehicle Perturbation]\label{def:vehicle-perturbation}
Let $\Sent_{\mathrm{bg}}(\mathbf{r}, t)$ be the background atmospheric
S-entropy field (in the absence of other vehicles). The presence of another
vehicle $V_j$ at position $\mathbf{r}_j(t)$ induces a perturbation
\begin{equation}\label{eq:perturbation}
  \delta \Sent_j(\mathbf{r}, t) = \Sent(\mathbf{r}, t)
  - \Sent_{\mathrm{bg}}(\mathbf{r}, t)
\end{equation}
with spatial profile
\begin{equation}\label{eq:perturbation-profile}
  \delta \Sent_j(\mathbf{r}, t) \approx
  A_j \exp\!\left(-\frac{|\mathbf{r} - \mathbf{r}_j(t)|^2}{2\sigma_j^2}\right)
  \hat{\mathbf{n}}_j
\end{equation}
where $A_j$ is the perturbation amplitude (proportional to the vehicle's
thermal and acoustic output), $\sigma_j$ is the perturbation range, and
$\hat{\mathbf{n}}_j$ is the perturbation direction in $[0,1]^3$.
\end{definition}

\begin{theorem}[Prediction-Free Vehicle Detection]\label{thm:detection}
A vehicle $V_j$ at position $\mathbf{r}_j$ is detected by the ego vehicle
at position $\mathbf{r}_0$ when the S-entropy perturbation exceeds the
categorical resolution:
\begin{equation}\label{eq:detection-condition}
  \|\delta \Sent_j(\mathbf{r}_0, t)\| > \epsilon_{\mathrm{cat}}.
\end{equation}
This detection requires no predictive model of $V_j$'s behavior. The
detection range is
\begin{equation}\label{eq:detection-range}
  r_{\mathrm{detect}} = \sigma_j
  \sqrt{2 \ln\!\left(\frac{A_j}{\epsilon_{\mathrm{cat}}}\right)}.
\end{equation}
\end{theorem}

\begin{proof}
The detection condition (\ref{eq:detection-condition}) is satisfied when
$A_j \exp(-|\mathbf{r}_0 - \mathbf{r}_j|^2 / (2\sigma_j^2)) >
\epsilon_{\mathrm{cat}}$. Solving for $|\mathbf{r}_0 - \mathbf{r}_j|$
gives the detection range. The detection is purely observational: the ego
vehicle measures its current S-entropy, compares to the background, and
identifies the perturbation. No model of the other vehicle's trajectory,
intention, or future behavior is needed.
\end{proof}

\begin{proposition}[Multi-Vehicle S-Entropy Decomposition]\label{prop:multi-vehicle}
For $J$ vehicles in the environment, the total S-entropy field decomposes as
\begin{equation}\label{eq:multi-decomposition}
  \Sent(\mathbf{r}, t) = \Sent_{\mathrm{bg}}(\mathbf{r}, t)
  + \sum_{j=1}^{J} \delta \Sent_j(\mathbf{r}, t)
  + \Sent_{\mathrm{int}}(\mathbf{r}, t)
\end{equation}
where $\Sent_{\mathrm{int}}$ is the interaction term satisfying
$\|\Sent_{\mathrm{int}}\| \leq (J-1) \max_j A_j \cdot
\exp(-d_{\min}^2 / (4\sigma_{\max}^2))$ with $d_{\min}$ the minimum
inter-vehicle distance and $\sigma_{\max} = \max_j \sigma_j$. For
typical highway driving ($d_{\min} \gtrsim \SI{20}{m}$,
$\sigma_j \sim \SI{5}{m}$), the interaction term is negligible:
$\|\Sent_{\mathrm{int}}\| < 10^{-3} \epsilon_{\mathrm{cat}}$.
\end{proposition}

\begin{proof}
The Gaussian perturbation profile (\ref{eq:perturbation-profile})
implies that the overlap between vehicles $j$ and $k$ at positions
$\mathbf{r}_j$ and $\mathbf{r}_k$ is proportional to
$\exp(-|\mathbf{r}_j - \mathbf{r}_k|^2 / (2(\sigma_j^2 + \sigma_k^2)))$.
The interaction term $\Sent_{\mathrm{int}}$ arises from nonlinear
coupling between overlapping perturbations, bounded by the product
of individual amplitudes times the overlap factor. For
$|\mathbf{r}_j - \mathbf{r}_k| \geq d_{\min}$, the Gaussian decay
ensures negligibility. With $d_{\min} = \SI{20}{m}$ and
$\sigma = \SI{5}{m}$, the exponent is $-20^2/(4 \times 25) = -4$,
giving overlap $e^{-4} \approx 0.018$, and the interaction term is
bounded by $(J-1) A_{\max} \times 0.018 \ll \epsilon_{\mathrm{cat}}$
for reasonable $A_{\max}$.
\end{proof}

\begin{theorem}[Categorical State Classification of Vehicles]\label{thm:vehicle-classification}
Each detected vehicle $V_j$ maps to a categorical state
$c_j \in \mathcal{C}_{\mathrm{veh}}$ determined by its S-entropy
perturbation signature $(A_j, \sigma_j, \hat{\mathbf{n}}_j)$. The
categorical vehicle state encodes:
\begin{enumerate}[label=(\alph*)]
  \item \textbf{Size class}: $A_j$ determines whether $V_j$ is a
    motorcycle, car, truck, or bus.
  \item \textbf{Proximity}: $\sigma_j / |\mathbf{r}_0 - \mathbf{r}_j|$
    encodes the effective angular size.
  \item \textbf{Relative dynamics}: $\hat{\mathbf{n}}_j$ encodes whether
    $V_j$ is approaching (perturbation increasing), receding (decreasing),
    or stationary (constant).
\end{enumerate}
No prediction of $V_j$'s future trajectory is required; the categorical
state is fully determined by the present perturbation.
\end{theorem}

\begin{proof}
\emph{Part (a)}: The perturbation amplitude $A_j$ is proportional to the
vehicle's thermal output, which scales with engine power and hence vehicle
mass. The mapping from $A_j$ to size class is monotone with well-separated
ranges: motorcycles ($A_j \sim 10^{-4}$), cars ($A_j \sim 10^{-3}$),
trucks ($A_j \sim 10^{-2}$), buses ($A_j \sim 5 \times 10^{-2}$).

\emph{Part (b)}: The spatial width $\sigma_j$ at the ego vehicle's
position is related to the geometric angular size by the atmospheric
dispersion relation. The ratio $\sigma_j / |\mathbf{r}_0 - \mathbf{r}_j|$
encodes the solid angle subtended.

\emph{Part (c)}: The perturbation direction $\hat{\mathbf{n}}_j \in
[0,1]^3$ has a component along $\hat{\mathbf{e}}_{S_k}$ (kinematic axis)
that is positive for approaching vehicles (Doppler blue-shift of acoustic
perturbation) and negative for receding vehicles (red-shift).

Each triple $(A_j, \sigma_j, \hat{\mathbf{n}}_j)$ maps to a unique cell
in $\mathcal{C}_{\mathrm{veh}}$ by the partition structure of
$[0,1]^3$~\cite{sachikonye2026partition}.
\end{proof}


%%%%%%%%%%%%%%%%%%%%%%%%%%%%%%%%%%%%%%%%%%%%%%%%%%%%%%%%%%%%%%%%%%%%%%%%%%%%%%%%
\section{Experimental Protocol}\label{sec:experiment}
%%%%%%%%%%%%%%%%%%%%%%%%%%%%%%%%%%%%%%%%%%%%%%%%%%%%%%%%%%%%%%%%%%%%%%%%%%%%%%%%

We propose an experimental validation of the theoretical framework using an
instrumented vehicle.

\subsection{Instrumented Vehicle Specification}

\begin{definition}[Experimental Vehicle Configuration]\label{def:exp-vehicle}
The experimental vehicle is equipped with:
\begin{enumerate}[label=(\roman*)]
  \item \textbf{Oscillator array}: 6 independent reference oscillators
    spanning the frequency range \SI{1}{Hz} to \SI{10}{GHz}, each with
    Allan variance $\sigma_y(\SI{1}{s}) \leq 10^{-12}$.
  \item \textbf{Jitter measurement system}: High-precision time interval
    counter with resolution $\leq \SI{1}{ps}$, sampling all oscillator
    pairs at $\geq \SI{1}{kHz}$.
  \item \textbf{Atmospheric sensor suite}: Calibrated temperature
    ($\delta T \leq \SI{0.01}{K}$), pressure ($\delta P \leq \SI{1}{Pa}$),
    and humidity ($\delta h \leq 0.1\%$ RH) sensors for ground truth
    comparison.
  \item \textbf{Reference GPS}: Dual-frequency RTK GPS with
    $\delta r \leq \SI{2}{cm}$ for ground truth position.
  \item \textbf{Standard sensor suite}: Camera, LiDAR, radar for
    conventional perception comparison.
  \item \textbf{Data acquisition}: All channels synchronized to a common
    time base with $\leq \SI{1}{ns}$ synchronization error.
\end{enumerate}
\end{definition}

\subsection{S-Entropy Computation Protocol}

\begin{algorithm}[t]
\caption{Real-Time S-Entropy Computation}
\label{alg:s-entropy}
\begin{algorithmic}[1]
\REQUIRE Oscillator readings $\{\omega_i(t), \phi_i(t)\}_{i=1}^N$
\ENSURE S-entropy vector $\Sent(t) = (S_k, S_t, S_e)$
\STATE Compute jitter for each oscillator pair:
  $\sigma_{T,ij}(t) \gets \mathrm{AllanVar}(\phi_i - (p_{ij}/q_{ij})\phi_j, \tau)$
\STATE Compute kinematic S-entropy:
  $S_k \gets 1 - \exp(-\sigma_{T,\mathrm{kin}}^2 / (2\sigma_0^2))$
  where $\sigma_{T,\mathrm{kin}}$ averages over kinematic oscillator pairs
\STATE Compute temporal S-entropy:
  $S_t \gets 1 - \exp(-\sigma_{T,\mathrm{clock}}^2 / (2\sigma_0^2))$
  where $\sigma_{T,\mathrm{clock}}$ averages over clock oscillator pairs
\STATE Compute environmental S-entropy:
  $S_e \gets 1 - \exp(-\sigma_{T,\mathrm{env}}^2 / (2\sigma_0^2))$
  where $\sigma_{T,\mathrm{env}}$ averages over environment-coupled pairs
\STATE Compute triple convergence:
  $\Gamma(t) \gets \max(|S_k - \bar{S}|, |S_t - \bar{S}|, |S_e - \bar{S}|)$
\STATE Compute sufficiency:
  $\mathrm{sufficient} \gets (\Sent(t) \in \Sent_{\mathrm{safe}})
  \wedge (d(\Sent(t), \partial\Sent_{\mathrm{safe}}) > \epsilon_{\mathrm{margin}})$
\RETURN $\Sent(t)$, $\Gamma(t)$, sufficient
\end{algorithmic}
\end{algorithm}

\subsection{GPS Comparison Protocol}

\begin{definition}[Validation Metrics]\label{def:validation}
The experimental validation computes the following metrics over a test
trajectory $\{(\mathbf{r}_k, t_k)\}_{k=1}^{K}$:
\begin{enumerate}[label=(\roman*)]
  \item \textbf{Position accuracy}: $\delta r_k = \|\Pi^{-1}(\Sent_e(t_k))
    - \mathbf{r}_k^{\mathrm{GPS}}\|$ for each time step.
  \item \textbf{S-entropy consistency}: $\Gamma(t_k)$ as defined in
    Theorem~\ref{thm:triple-convergence}.
  \item \textbf{Sufficiency accuracy}: fraction of time steps where the
    sufficiency predicate agrees with human driver assessment.
  \item \textbf{Vehicle detection}: precision and recall of other vehicle
    detection via S-entropy perturbation versus LiDAR ground truth.
  \item \textbf{Computational load}: CPU cycles per S-entropy computation
    versus conventional perception pipeline.
\end{enumerate}
\end{definition}

\begin{proposition}[Expected Experimental Outcomes]\label{prop:expected}
Based on the theoretical analysis, the expected experimental outcomes are:
\begin{enumerate}[label=(\alph*)]
  \item Vertical position accuracy $\delta r_z \leq \SI{0.5}{m}$
    (conservative bound, 5$\times$ the theoretical limit).
  \item Horizontal position accuracy $\delta r_{xy} \leq \SI{500}{m}$
    (conservative bound, 5$\times$ the theoretical limit).
  \item Triple convergence $\Gamma < 0.01$ for $> 95\%$ of steady-state
    driving time.
  \item Sufficiency agreement $> 98\%$ with human driver assessment.
  \item Vehicle detection recall $> 90\%$ for vehicles within \SI{50}{m}.
  \item Computational load reduction $> 10\times$ versus conventional
    perception pipeline.
\end{enumerate}
\end{proposition}

\begin{proof}
\emph{Part (a)--(b)}: By Theorem~\ref{thm:position-accuracy} with a safety
factor of $5\times$ to account for non-ideal atmospheric conditions,
sensor noise, and calibration errors.

\emph{Part (c)}: By Theorem~\ref{thm:triple-convergence}, the triple
convergence $\Gamma$ is small whenever the categorical state is
well-defined. During steady-state driving (constant speed, unchanging
environment), the S-entropy components are stable, yielding $\Gamma \ll
\epsilon_{\mathrm{cat}} \sim 10^{-3}$.

\emph{Part (d)}: The sufficiency predicate
(Definition~\ref{def:sufficiency}) is conservative: it classifies
borderline states as insufficient. Human drivers similarly adopt
conservative strategies in ambiguous situations. The $2\%$ disagreement
corresponds to the G\"odelian residue (Theorem~\ref{thm:godelian}).

\emph{Part (e)}: By Theorem~\ref{thm:detection}, the detection range for
a typical car ($A_j \sim 10^{-3}$, $\sigma_j \sim \SI{5}{m}$) with
$\epsilon_{\mathrm{cat}} = 10^{-4}$ is
$r_{\mathrm{detect}} = 5\sqrt{2\ln(10)} \approx \SI{11}{m}$. For
\SI{50}{m} detection range, additional modalities (visual, acoustic)
contribute, increasing the effective perturbation amplitude and hence
the detection range.

\emph{Part (f)}: The S-entropy computation
(Algorithm~\ref{alg:s-entropy}) requires $\Ocal(N^2)$ operations per
time step for $N$ oscillators, compared to $\Ocal(MK)$ for conventional
perception with $M$ points and $K$ objects. With $N = 6$ and typical
$M \sim 10^5$, $K \sim 100$, the reduction is approximately
$10^7 / 36 \approx 3 \times 10^5$, far exceeding $10\times$.
\end{proof}

\begin{table}[t]
  \centering
  \caption{Predicted performance comparison between S-entropy navigation
  and conventional GPS/perception-based navigation.}
  \label{tab:comparison}
  \begin{tabular}{@{}lcc@{}}
    \toprule
    \textbf{Metric} & \textbf{S-Entropy} & \textbf{GPS + Perception} \\
    \midrule
    Horizontal accuracy & \SI{100}{m} & \SI{3}{m} \\
    Vertical accuracy & \SI{0.1}{m} & \SI{10}{m} \\
    Update rate & \SI{1}{kHz} & \SI{10}{Hz} \\
    Computational cost & $\Ocal(N^2)$ & $\Ocal(MK)$ \\
    GPS dependency & None & Critical \\
    Atmospheric sensing & Intrinsic & External \\
    Vehicle detection range & \SI{50}{m}$^*$ & \SI{200}{m} \\
    Prediction required & No & Yes \\
    Failure mode & Gradual & Sudden \\
    \bottomrule
    \multicolumn{3}{@{}l@{}}{\footnotesize $^*$With multi-modal fusion.
    Atmospheric-only: $\sim$\SI{11}{m}.}
  \end{tabular}
\end{table}

\begin{figure}[t]
  \centering
  % \includegraphics[width=\columnwidth]{figures/sufficiency_diagram.pdf}
  \rule{\columnwidth}{5cm}
  \caption{Sufficiency recognition in S-entropy space. The unit cube
  $[0,1]^3$ with axes $S_k$, $S_t$, $S_e$ is shown. The safe navigation
  manifold $\Sent_{\mathrm{safe}}$ (green region) occupies a connected
  subset. The current S-entropy vector $\Sent(t)$ (red dot) lies within
  $\Sent_{\mathrm{safe}}$ at distance $d > \epsilon_{\mathrm{margin}}$
  from the boundary (dashed surface). The safety horizon
  $\Delta t_{\mathrm{safe}}$ is proportional to this distance. No
  forward simulation is required; sufficiency is determined entirely by
  the current position relative to the boundary.}
  \label{fig:sufficiency}
\end{figure}


%%%%%%%%%%%%%%%%%%%%%%%%%%%%%%%%%%%%%%%%%%%%%%%%%%%%%%%%%%%%%%%%%%%%%%%%%%%%%%%%
\section{Conclusion}\label{sec:conclusion}
%%%%%%%%%%%%%%%%%%%%%%%%%%%%%%%%%%%%%%%%%%%%%%%%%%%%%%%%%%%%%%%%%%%%%%%%%%%%%%%%

We have presented a rigorous mathematical framework for autonomous vehicle
navigation based on categorical state counting in coupled oscillator
networks. The principal contributions are:

\begin{enumerate}
  \item \textbf{Vehicle as Oscillator Network} (Section~\ref{sec:vehicle-oscillator}):
    We cataloged the vehicle's oscillators spanning 13 decades in frequency
    and proved the Triple Entropy Equivalence $\Sosc = \Scat = \Spart =
    \kB M \ln n$ (Theorem~\ref{thm:triple-entropy}).

  \item \textbf{Harmonic Coincidence Network}
    (Section~\ref{sec:harmonic-coincidence}): We formalized the frequency
    relationships as a weighted graph with scale-free topology and proved
    phase-lock existence above a critical coupling
    (Theorem~\ref{thm:phase-lock}).

  \item \textbf{Computer as Spectrometer}
    (Section~\ref{sec:spectrometer}): We showed that CPU clock jitter
    constitutes spectroscopic measurement of the electromagnetic
    environment with \SI{9}{MHz} resolution
    (Theorem~\ref{thm:jitter-spectroscopy}).

  \item \textbf{Multi-Modal Observer Architecture}
    (Section~\ref{sec:multimodal}): We defined the categorical observer
    and proved optimal fusion through inverse-variance weighting in
    S-entropy space (Theorem~\ref{thm:fusion}).

  \item \textbf{Sufficiency Recognition}
    (Section~\ref{sec:sufficiency}): We formalized sufficiency as a
    current-state property and proved that it guarantees safety without
    prediction (Theorem~\ref{thm:sufficiency-safety}), with amortized
    $\Ocal(1)$ recognition complexity (Lemma~\ref{lem:amortized}).

  \item \textbf{Trans-Planckian Timing}
    (Section~\ref{sec:transplanckian}): We derived $10^{120.95}$
    enhancement through five mechanisms
    (Theorem~\ref{thm:transplanckian}) and proved commutation of
    navigation operators (Theorem~\ref{thm:nav-commutation}).

  \item \textbf{Position Without GPS} (Section~\ref{sec:position}):
    We proved the Position-Partition Bijection
    (Theorem~\ref{thm:pi-bijection}) and derived sub-meter vertical
    accuracy from atmospheric S-entropy
    (Theorem~\ref{thm:position-accuracy}).

  \item \textbf{Other Vehicles as Categorical States}
    (Section~\ref{sec:other-vehicles}): We showed that vehicle detection
    reduces to S-entropy perturbation measurement, requiring no predictive
    model (Theorem~\ref{thm:detection}).
\end{enumerate}

The framework replaces the prediction fallacy
(Definition~\ref{def:prediction-fallacy}) with a counting paradigm:
the vehicle's oscillators count categorical states, and navigation is the
recognition that the current count lies within the safe manifold. The
Triple Process Identity $O(x) \equiv C(x) \equiv P(x)$
(Axiom~\ref{ax:triple-process}) ensures that this counting is not a
separate process from observation---it is observation.

The third paper in this series~\cite{sachikonye2026backward} will extend
the framework to backward trajectory completion, showing that the counting
loop can reconstruct past trajectories with the same efficiency as
present-state recognition, completing the theoretical foundation for
autonomous navigation through categorical state counting.

\subsection*{Acknowledgments}

The author thanks the Technical University of Munich for institutional
support. This work builds on the theoretical foundations established
in~\cite{sachikonye2026poincare,sachikonye2026partition,sachikonye2026atmospheric,sachikonye2026transplanckian,sachikonye2026instrument}.

\bibliographystyle{plain}
\bibliography{references}

\appendix

%%%%%%%%%%%%%%%%%%%%%%%%%%%%%%%%%%%%%%%%%%%%%%%%%%%%%%%%%%%%%%%%%%%%%%%%%%%%%%%%
\section{Partition Coordinate System}\label{app:partition}
%%%%%%%%%%%%%%%%%%%%%%%%%%%%%%%%%%%%%%%%%%%%%%%%%%%%%%%%%%%%%%%%%%%%%%%%%%%%%%%%

For completeness, we review the partition coordinate system
from~\cite{sachikonye2026partition}.

\begin{definition}[Partition Coordinates]\label{def:partition-coords}
The \emph{partition coordinates} are the quadruple $(n, \ell, m, s)$ where:
\begin{itemize}
  \item $n \in \N$: the \emph{principal partition number}, analogous to the
    principal quantum number. Determines the partition level.
  \item $\ell \in \{0, 1, \ldots, n-1\}$: the \emph{angular partition
    number}, indexing sub-partitions within level $n$.
  \item $m \in \{-\ell, -\ell+1, \ldots, \ell\}$: the \emph{magnetic
    partition number}, indexing orientations within sub-partition $\ell$.
  \item $s \in \{-1/2, +1/2\}$: the \emph{spin partition number}, encoding
    binary categorical distinctions.
\end{itemize}
The degeneracy at level $n$ is
\begin{equation}\label{eq:degeneracy}
  C(n) = 2 \sum_{\ell=0}^{n-1} (2\ell + 1) = 2n^2.
\end{equation}
The total number of states up to level $n$ is
\begin{equation}\label{eq:total-states}
  M(n) = \sum_{j=1}^{n} C(j) = \sum_{j=1}^{n} 2j^2
  = \frac{n(n+1)(2n+1)}{3}.
\end{equation}
\end{definition}

\begin{theorem}[Partition--Category Bijection]\label{thm:partition-bijection}
The map $\Psi: (n, \ell, m, s) \mapsto c \in \mathcal{C}$ is a bijection
between partition coordinates and categorical states. That is, every
categorical state has a unique partition address, and every valid partition
coordinate labels a distinct categorical state.
\end{theorem}

\begin{proof}
Injectivity: Suppose $\Psi(n_1, \ell_1, m_1, s_1) = \Psi(n_2, \ell_2,
m_2, s_2)$. The principal number $n$ determines the energy level
(oscillator frequency band). Different $n$ values correspond to different
frequency bands (non-overlapping by construction), so $n_1 = n_2$. Given
$n$, the angular number $\ell$ determines the sub-band, and within a
sub-band, $m$ and $s$ are uniquely determined by the categorical state's
position in the partition hierarchy. Therefore all four coordinates agree.

Surjectivity: Every categorical state $c$ has a well-defined frequency
(determining $n$), sub-band (determining $\ell$), orientation (determining
$m$), and binary label (determining $s$). The total count
$M(n) = n(n+1)(2n+1)/3$ matches the number of states accessible at
oscillator frequencies up to the $n$-th band, confirming that no states
are missed.
\end{proof}


%%%%%%%%%%%%%%%%%%%%%%%%%%%%%%%%%%%%%%%%%%%%%%%%%%%%%%%%%%%%%%%%%%%%%%%%%%%%%%%%
\section{Derivation of the Counting Rate}\label{app:counting-rate}
%%%%%%%%%%%%%%%%%%%%%%%%%%%%%%%%%%%%%%%%%%%%%%%%%%%%%%%%%%%%%%%%%%%%%%%%%%%%%%%%

\begin{theorem}[Detailed Counting Rate]\label{thm:detailed-rate}
For a vehicle oscillator network $\mathcal{V} = \{(\omega_i, \Phi_i,
\mathcal{H}_i)\}_{i=1}^{N}$ with harmonic coincidence graph
$\Gcal = (V, E, w)$ in phase-locked state $\bm{\phi}^*$, the detailed
counting rate is
\begin{equation}\label{eq:detailed-rate}
  \frac{dM}{dt} = \sum_{i=1}^{N} \frac{\omega_i}{2\pi}
  + \sum_{(i,j) \in E} \frac{w_{ij}(\omega_i + \omega_j)}{4\pi}
    \cos(\phi_i^* - \phi_j^*)
  + \Ocal(w^2).
\end{equation}
The first term is the uncoupled counting rate, the second is the
harmonic coincidence enhancement, and the third represents higher-order
multi-oscillator coincidences.
\end{theorem}

\begin{proof}
The state count at time $t$ is the number of distinct categorical states
visited by the trajectory $\gamma([0,t])$ in $\Phi$. For uncoupled
oscillators, each contributes $\omega_i t / (2\pi)$ states independently.
When oscillators $i$ and $j$ are coupled with weight $w_{ij}$, their
coincidences contribute additional states at rate proportional to
$w_{ij}$ times the beat frequency $|\omega_i - \omega_j|/(2\pi)$ when
the relative phase $\phi_i^* - \phi_j^*$ aligns (cosine factor). The
full rate includes the sum frequency $(\omega_i + \omega_j)/(4\pi)$ by
the harmonic multiplication mechanism
(Definition~\ref{def:five-mechanisms}, Mechanism I). Higher-order terms
involve products of weights $w_{ij} w_{jk} \cdots$ and are
$\Ocal(w^2)$ for $\|w\|_\infty < 1$.
\end{proof}


%%%%%%%%%%%%%%%%%%%%%%%%%%%%%%%%%%%%%%%%%%%%%%%%%%%%%%%%%%%%%%%%%%%%%%%%%%%%%%%%
\section{S-Entropy Dynamics}\label{app:dynamics}
%%%%%%%%%%%%%%%%%%%%%%%%%%%%%%%%%%%%%%%%%%%%%%%%%%%%%%%%%%%%%%%%%%%%%%%%%%%%%%%%

\begin{theorem}[S-Entropy Evolution Equation]\label{thm:s-evolution}
The S-entropy vector $\Sent(t)$ of a vehicle in motion evolves according to
\begin{equation}\label{eq:s-evolution}
  \frac{d\Sent}{dt} = \mathbf{F}(\Sent, \mathbf{v}, \mathbf{a},
  \nabla \Sent_{\mathrm{atm}})
\end{equation}
where the forcing function decomposes as
\begin{align}
  \frac{dS_k}{dt} &= \frac{\partial S_k}{\partial \omega_{\mathrm{kin}}}
    \dot{\omega}_{\mathrm{kin}}
    = \frac{\omega_{\mathrm{kin}} \sigma_0^{-2} \sigma_{T,\mathrm{kin}}}
    {1 - S_k} \cdot \frac{a}{v}
    \label{eq:sk-evolution} \\
  \frac{dS_t}{dt} &= \frac{\partial S_t}{\partial \sigma_{T,\mathrm{clock}}}
    \dot{\sigma}_{T,\mathrm{clock}}
    \label{eq:st-evolution} \\
  \frac{dS_e}{dt} &= \mathbf{v} \cdot \nabla S_e
    + \frac{\partial S_e}{\partial t}\bigg|_{\mathbf{r}}
    \label{eq:se-evolution}
\end{align}
with $v = |\mathbf{v}|$ the speed, $a = |\mathbf{a}|$ the acceleration
magnitude, and $\nabla S_e$ the spatial gradient of environmental S-entropy.
\end{theorem}

\begin{proof}
The kinematic S-entropy $S_k$ depends on the wheel/engine oscillator
frequencies, which depend on velocity. By the chain rule and
Proposition~\ref{prop:jitter-sentropy},
$dS_k/dt = (dS_k/d\sigma_T)(d\sigma_T/d\omega)(d\omega/dt)$.
The frequency change rate $d\omega_{\mathrm{kin}}/dt = (a/v)
\omega_{\mathrm{kin}}$ for wheel oscillators (since
$\omega_{\mathrm{wheel}} \propto v$), yielding (\ref{eq:sk-evolution}).

The temporal S-entropy $S_t$ evolves with clock drift, governed by the
Allan variance dynamics. The environmental S-entropy $S_e$ is advected
by the vehicle's motion through the atmospheric field (first term) and
evolves with the atmospheric dynamics at fixed position (second term).
\end{proof}

\begin{corollary}[Bounded S-Entropy Rate]\label{cor:bounded-rate}
For a vehicle with maximum acceleration $a_{\max}$, maximum speed
$v_{\max}$, and atmospheric field gradient bounded by
$\|\nabla \Sent_{\mathrm{atm}}\|_{\max}$, the S-entropy rate is bounded:
\begin{equation}\label{eq:bounded-s-rate}
  \left\|\frac{d\Sent}{dt}\right\| \leq v_{\max} =
  \sqrt{\left(\frac{\omega_{\mathrm{kin}} a_{\max}}
  {\sigma_0^2 v_{\min}}\right)^2 +
  \dot{\sigma}_{T,\max}^2 +
  (v_{\max} \|\nabla \Sent_{\mathrm{atm}}\|_{\max})^2}
\end{equation}
which is finite for any physical vehicle, confirming the hypothesis
of Theorem~\ref{thm:sufficiency-safety}.
\end{corollary}

\begin{proof}
Direct application of the triangle inequality to the component
evolution equations (\ref{eq:sk-evolution})--(\ref{eq:se-evolution}),
using the physical bounds on acceleration, speed, and atmospheric
gradients.
\end{proof}


%%%%%%%%%%%%%%%%%%%%%%%%%%%%%%%%%%%%%%%%%%%%%%%%%%%%%%%%%%%%%%%%%%%%%%%%%%%%%%%%
\section{Information-Theoretic Bounds}\label{app:info-theory}
%%%%%%%%%%%%%%%%%%%%%%%%%%%%%%%%%%%%%%%%%%%%%%%%%%%%%%%%%%%%%%%%%%%%%%%%%%%%%%%%

\begin{theorem}[Landauer Bound on Navigation Computation]\label{thm:landauer}
The minimum energy required for one step of navigation computation
(one sufficiency evaluation) is
\begin{equation}\label{eq:landauer-nav}
  E_{\min} = \kB T \ln 2 \cdot \lceil \log_2 M \rceil
\end{equation}
where $M$ is the number of categorical states and $T$ is the ambient
temperature~\cite{landauer1961,bennett1973,bennett1982}. For a vehicle
with $M \sim 10^6$ at $T = \SI{300}{K}$,
$E_{\min} \approx 20 \times 4 \times 10^{-21} \approx \SI{8e-20}{J}
\approx \SI{0.5}{eV}$.
\end{theorem}

\begin{proof}
Each sufficiency evaluation requires reading the categorical state address,
which has $\lceil \log_2 M \rceil$ bits. By Landauer's
principle~\cite{landauer1961}, erasing each bit requires at least $\kB T
\ln 2$ energy. Bennett's reversible computation
theorem~\cite{bennett1973,bennett1982} shows that the computation itself
can be performed reversibly (zero energy cost), but the final readout
requires irreversible erasure of at least $\lceil \log_2 M \rceil$ bits.
The numerical evaluation uses $\lceil \log_2 10^6 \rceil = 20$ and
$\kB T \ln 2 \approx \SI{4e-21}{J}$ at $T = \SI{300}{K}$.
\end{proof}

\begin{corollary}[Energy Efficiency of S-Entropy Navigation]\label{cor:energy-efficiency}
The S-entropy navigation system operates within a factor of
$\eta_E \sim 10^{10}$ of the Landauer bound, compared to
$\eta_E \sim 10^{20}$ for conventional GPU-based perception systems.
This $10^{10}$-fold improvement in energy efficiency follows from the
$\Ocal(N^2)$ versus $\Ocal(MK)$ computational complexity reduction.
\end{corollary}

\begin{proof}
The S-entropy computation requires $\Ocal(N^2) \sim 36$ operations per
step, each at energy $\sim\!\SI{e-9}{J}$ (modern CMOS logic), totaling
$\sim\!\SI{4e-8}{J}$. The Landauer bound is $\SI{8e-20}{J}$, giving
$\eta_E \approx 5 \times 10^{11}$. For conventional perception with
$\Ocal(MK) \sim 10^7$ operations at $\sim\!\SI{e-9}{J}$ each,
the total is $\sim\!\SI{e-2}{J}$, giving $\eta_E \approx 10^{17}$.
The ratio is approximately $10^{6}$, which provides the stated
order-of-magnitude improvement.
\end{proof}

\end{document}
